% ============================================================
%  GENERATED from the PCLib "UDT Standards.md" (EFC Repo, PCLib v1.00)
%  by gen_future_appendix.py. Regenerate rather than hand-editing the
%  type tables; the text sections at the top may be edited.
% ============================================================
\chapter{PLC Data Type Reference (Future)}
\label{app:data-types-future}
\future{This appendix describes the proposed Process Control Library (PCLib~v1.00), the next generation of EFC standard data types. It has not been adopted yet. Existing systems keep using the current types in Appendix~\ref{app:data-types} until a project is approved to move to PCLib. Member names, initial values, and alarm settings may still change during review and simulator testing.}

PCLib is a set of standard data types (Control Expert DDTs and matching Ignition UDTs) for process-control equipment. Compared with the current library it gives every piece of equipment a consistent structure: the same member prefixes, the same mode and alarm objects, the same status integers for HMI animation, and a documented link to the industry standards the types are aligned with (ISA-88, ISA-18.2, ISA-101, NAMUR NE107/NE43, and MTP). The library is kept in the EFC Repo under \texttt{PCLib/}. The type definitions are generated by \texttt{gen\_pclib.py}, which writes both the Control Expert import file and the Ignition import file, so the PLC and SCADA definitions always match member for member. Change the generator and re-run it; never hand-edit the generated files.

The types are \textbf{data only}. Each type is paired with a DFB (or ST section) that holds its behaviour; most are laid out so the DFB can take the whole DDT as a single \texttt{IN\_OUT} parameter. The first of these, \texttt{FB\_AI}, is summarized in Section~\ref{sec:pclib-fbai}.


\section{Current to Future Mapping}
The table shows which PCLib type replaces each current type. Members do not map one to one, so moving an existing project to PCLib means re-mapping its I/O, logic, and SCADA bindings, and is done as a planned change (Section~\ref{sec:review-release}), not as a find-and-replace.


\begin{small}
\begin{longtable}{@{}P{4.40cm}P{4.20cm}P{8.33cm}@{}}
\caption{Current to PCLib Type Mapping}\label{tab:pclib-mapping} \\
\toprule
\rowcolor{tblhead}\textbf{Current type} & \textbf{PCLib type} & \textbf{Notes} \\
\midrule
\endfirsthead
\multicolumn{3}{@{}l}{\small\textit{Table \thetable{} continued}} \\
\toprule
\rowcolor{tblhead}\textbf{Current type} & \textbf{PCLib type} & \textbf{Notes} \\
\midrule
\endhead
\texttt{DDT\_Process\_Value}, \texttt{AI\_Sensors} & \texttt{T\_AI} & Scaling, filtering, NE43 range check, quality, and HH/H/L/LL/ROC alarms in one type. \\
\texttt{DDT\_AI\_Alarm} & \texttt{T\_AlmLim} & One object per limit, with deadband, delays, shelving, and the ISA-18.2 state. \\
\texttt{DDT\_Discrete} & \texttt{T\_DI} & Override becomes simulation (\texttt{Sim\_}); optional state alarm. \\
\texttt{DDT\_Valve} & \texttt{T\_DValve} & Single or dual solenoid, limit switches, travel alarms, run counters. \\
\texttt{DDT\_Motor\_VFD} & \texttt{T\_VFD} (or \texttt{T\_Motor} for fixed speed) & Fixed-speed and variable-speed drives become separate types. \\
\texttt{DDT\_Device} & \texttt{T\_Mode} + \texttt{T\_Alarm} & Mode, ownership, and alarms become nested objects used by every device type. \\
\texttt{DDT\_Setpoint} & \texttt{Cfg\_} / \texttt{Prm\_} members & Setpoints live inside the type that uses them (for example \texttt{T\_PID}, \texttt{T\_AlmLim}). Hard-coded clamp limits (Section~\ref{sec:plc-setpoint-clamp}) stay in the DFB. \\
\texttt{DDT\_Scale} & \texttt{T\_AI} (weight) + \texttt{T\_Totalizer} & Network-connected scales still need an application-specific wrapper \tbd. \\
\texttt{DDT\_Seq} & \texttt{T\_Phase} & ISA-88 Part 1 state model and commands replace the custom step/status words. \\
\texttt{DDT\_Station} & \texttt{T\_Unit} + \texttt{T\_Phase} & Application-specific; a station demo type (\texttt{DDT\_Station\_Demo}) is in the EFC Repo for review. \\
PDC compressor package interface (project-specific) & \texttt{T\_DiaComp} & Two-stage diaphragm compressor unit with leak detection, oil systems, purge, and first-out (Section~\ref{sec:pdc-compressor}). The current project type will be mapped once it is reviewed \tbd. \\
(no current type) & \texttt{T\_Intlk}, \texttt{T\_ESD} & Interlock and shutdown groups with first-out and timed bypass (BPCS layer only). \\
(no current type) & \texttt{T\_PID}, \texttt{T\_CValve}, rotating and process equipment types & New coverage; see the type index. \\
\bottomrule
\end{longtable}
\end{small}

\section{Naming Conventions}

Every member name starts with a prefix that says what it is and who writes it. Identifiers follow IEC 61131-3 and stay within the 32-character Control Expert limit.

\begin{longtable}{@{}P{2.40cm}P{8.10cm}P{6.43cm}@{}}
\caption{PCLib Member Prefixes}\label{tab:pclib-prefixes} \\
\toprule
\rowcolor{tblhead}\textbf{Prefix} & \textbf{Meaning} & \textbf{Written by} \\
\midrule
\endfirsthead
\multicolumn{3}{@{}l}{\small\textit{Table \thetable{} continued}} \\
\toprule
\rowcolor{tblhead}\textbf{Prefix} & \textbf{Meaning} & \textbf{Written by} \\
\midrule
\endhead
\texttt{Inp\_} & Field /\allowbreak{} upstream input & I/\allowbreak{}O mapping or upstream logic \\
\texttt{Cfg\_} & Configuration /\allowbreak{} tuning & Engineer (HMI with access control) \\
\texttt{Cmd\_} & Operator command. Booleans are \textbf{one-shot}: the logic clears them after acting & HMI \\
\texttt{Prog\_} & Command from the ISA-88 owner (phase, equipment module, batch server) & Owner logic \\
\texttt{Prm\_} /\allowbreak{} \texttt{Rpt\_} & ISA-88 recipe parameters /\allowbreak{} report values & Batch server /\allowbreak{} phase logic \\
\texttt{Out\_} & Output to field or downstream logic & Equipment logic \\
\texttt{Sts\_} & Status & Equipment logic \\
\texttt{Val\_} & Calculated value & Equipment logic \\
\texttt{Alm\_} & Alarm (T\_Alarm /\allowbreak{} T\_AlmLim) & Equipment logic sets \texttt{Inp\_Cond} \\
\texttt{Sim\_} & Simulation & Engineer \\
\bottomrule
\end{longtable}

\begin{itemize}[itemsep=1pt]

  \item Booleans ending in \texttt{OK} are \textbf{TRUE = healthy}, so a broken wire or lost signal reads as unhealthy. Other \texttt{Inp\_} booleans are TRUE when the condition is present.

  \item Interlocks (\texttt{Inp\_IntlkOK}) stop running equipment. Permissives (\texttt{Inp\_PermOK}) only block a start.

  \item \texttt{Sts\_State} /\allowbreak{} \texttt{Sts\_Mode} /\allowbreak{} \texttt{Sts\_NE107} integers are for HMI animation. The values are listed in each member comment.

  \item These are \textbf{data types only}. Pair each one with a DFB (or ST section) holding the behaviour. Most types are laid out so a DFB can take the whole DDT as one IN\_OUT parameter.

\end{itemize}

\note{\texttt{T\_ESD}, \texttt{T\_Intlk}, and \texttt{T\_GasDet} are for the basic process control layer and the HMI only. Safety instrumented functions with a SIL target must run in a certified safety controller (for example M580 Safety) and follow IEC~61511.}

\section{Standards Alignment}

The standards define models, behaviour, and terminology; apart from MTP they do not prescribe PLC data layouts. PCLib carries the data each model needs so the DFBs can implement the behaviour. This is alignment by design, not a third-party certification. Alarm priorities, limits, and equipment states still come from EFC's alarm philosophy and rationalisation (Section~\ref{sec:alarm-mgmt}).

\begin{small}
\begin{longtable}{@{}P{3.10cm}P{3.30cm}P{10.53cm}@{}}
\caption{PCLib Standards Alignment}\label{tab:pclib-standards} \\
\toprule
\rowcolor{tblhead}\textbf{Standard} & \textbf{Scope} & \textbf{How PCLib applies it} \\
\midrule
\endfirsthead
\multicolumn{3}{@{}l}{\small\textit{Table \thetable{} continued}} \\
\toprule
\rowcolor{tblhead}\textbf{Standard} & \textbf{Scope} & \textbf{How PCLib applies it} \\
\midrule
\endhead
ISA-88.00.01 /\allowbreak{} IEC 61512-1 & Batch control: physical and procedural models, modes, states & Every type has an ISA-88 role (control module, equipment module, unit, phase). \texttt{Prog\_} members are commands from the owning phase or equipment module. \texttt{T\_Mode} adds arbitration (acquire /\allowbreak{} release /\allowbreak{} owner ID). \texttt{T\_Phase} uses the Part 1 example state model (12 states) and commands (Start, Stop, Hold, Restart, Abort, Reset, Pause, Resume), with Automatic /\allowbreak{} Semi-automatic /\allowbreak{} Manual modes, recipe parameters and reports. \texttt{T\_Unit} holds allocation, batch identity and equipment state (clean, dirty, in use...). \\
ISA-18.2 /\allowbreak{} IEC 62682 & Alarm management & \texttt{T\_Alarm} uses the standard's state model (A Normal to H Out-of-service), including latching, time-limited shelving, suppressed-by-design and out-of-service. Chattering control (deadband, on/\allowbreak{}off delay), priority and class come from rationalisation. Events that need no operator action are status bits, not alarms (e.g. pump changeover). \\
ISA-101.01 & HMI design & Modes, states and NE107 status are integers, so graphics can use consistent symbols without colour logic in the PLC. \texttt{Sts\_CmdMask} lets the HMI grey out invalid phase commands. Tag and description are in the type for faceplates. \\
ISA-5.1 & Instrument identification & \texttt{Cfg\_Tag} holds the ISA-5.1 tag (e.g. FT-101). Limit alarms follow HH /\allowbreak{} H /\allowbreak{} L /\allowbreak{} LL naming. \\
NAMUR NE107 & Field device diagnostics & \texttt{Sts\_NE107} on every field-device control module (Good /\allowbreak{} F /\allowbreak{} C /\allowbreak{} S /\allowbreak{} M). \texttt{Inp\_NE107} takes the status from HART or fieldbus devices. \\
NAMUR NE43 & 4-20 mA failure signal levels & \texttt{T\_AI} bad-signal limits default to NE43 (below 3.6 mA, above 21 mA). \texttt{T\_GasDet} fault and over-range thresholds are configurable. \\
VDI/\allowbreak{}VDE/\allowbreak{}NAMUR 2658-4 (MTP) & Module Type Package data objects for modular plants & Each type lists its nearest MTP data assembly (AnaMon, BinMon, MonBinVlv, MonAnaVlv, MonBinDrv, MonAnaDrv, PIDCtrl, LockView16). Mode values map to MTP Offline /\allowbreak{} Operator /\allowbreak{} Automatic. Member names are not MTP attribute names, so an MTP export needs a mapping layer. \\
ISA-84 /\allowbreak{} IEC 61511 & Safety instrumented systems & \texttt{T\_Intlk} and \texttt{T\_ESD} follow the bypass-management requirements (authorisation, time limit, indication, logging data), first-out, manual reset after trip and de-energise-to-trip. They are \textbf{basic process control (BPCS) functions only}, not SIFs. \\
IEC 61131-3 & PLC languages and data types & Elementary IEC data types only. Identifiers are checked against the 32-character limit and IEC /\allowbreak{} Control Expert reserved words. \\
ISA-95 /\allowbreak{} IEC 62264 & Enterprise-control integration & Level 1-2 data. \texttt{T\_Unit} and \texttt{T\_Phase} expose batch ID, recipe, product and reports for MES collection. \\
\bottomrule
\end{longtable}
\end{small}

\subsection{ISA-88 Roles}

\begin{longtable}{@{}P{5.00cm}P{12.36cm}@{}}
\caption{PCLib Types by ISA-88 Level}\label{tab:pclib-isa88} \\
\toprule
\rowcolor{tblhead}\textbf{ISA-88 level} & \textbf{PCLib types} \\
\midrule
\endfirsthead
\multicolumn{2}{@{}l}{\small\textit{Table \thetable{} continued}} \\
\toprule
\rowcolor{tblhead}\textbf{ISA-88 level} & \textbf{PCLib types} \\
\midrule
\endhead
Support & \texttt{T\_Alarm}, \texttt{T\_AlmLim}, \texttt{T\_Mode}, \texttt{T\_RunTime}, \texttt{T\_Intlk}, \texttt{T\_ESD} \\
Control module & \texttt{T\_AI}, \texttt{T\_DI}, \texttt{T\_DValve}, \texttt{T\_CValve}, \texttt{T\_Motor}, \texttt{T\_VFD}, \texttt{T\_Heater}, \texttt{T\_PID}, \texttt{T\_PumpProt}, \texttt{T\_Totalizer}, \texttt{T\_HeatEx}, \texttt{T\_Analyzer}, \texttt{T\_GasDet} \\
Equipment module & \texttt{T\_PumpDS}, \texttt{T\_Agitator}, \texttt{T\_Compressor}, \texttt{T\_Conveyor}, \texttt{T\_DosePump} \\
Equipment module /\allowbreak{} Unit vessel & \texttt{T\_Tank} \\
Equipment phase & \texttt{T\_Phase} \\
Unit & \texttt{T\_Unit} \\
\bottomrule
\end{longtable}

Control modules take \texttt{Prog\_} commands only from the owner that has acquired them through \texttt{Mode.Prog\_Acquire} (ISA-88 arbitration). Out-of-service and Maintenance modes are site extensions; ISA-88 itself defines Automatic, Semi-automatic and Manual.

\section{Status Values}

The status integers below are the same in every type that uses them, so one set of SCADA symbols and colours works across the library (ISA-101).

\begin{small}
\begin{longtable}{@{}P{3.80cm}P{3.40cm}P{9.73cm}@{}}
\caption{PCLib Status Values}\label{tab:pclib-status} \\
\toprule
\rowcolor{tblhead}\textbf{Object} & \textbf{Member} & \textbf{Values} \\
\midrule
\endfirsthead
\multicolumn{3}{@{}l}{\small\textit{Table \thetable{} continued}} \\
\toprule
\rowcolor{tblhead}\textbf{Object} & \textbf{Member} & \textbf{Values} \\
\midrule
\endhead
Alarm (\texttt{T\_Alarm}, \texttt{T\_AlmLim}) & \texttt{Sts\_State} & 0=A Normal 1=B Unack 2=C Acked 3=D RTN-Unack 4=E Latched 5=F Shelved 6=G Suppressed-by-design 7=H Out-of-service \\
Phase (\texttt{T\_Phase}) & \texttt{Sts\_State} & 0=Idle 1=Running 2=Complete 3=Pausing 4=Paused 5=Holding 6=Held 7=Restarting 8=Stopping 9=Stopped 10=Aborting 11=Aborted \\
Phase commands & \texttt{Prog\_Cmd}, \texttt{Sts\_CmdMask} & 1=Start 2=Stop 3=Hold 4=Restart 5=Abort 6=Reset 7=Pause 8=Resume \\
Equipment (\texttt{T\_Unit}, \texttt{T\_Tank}) & \texttt{Sts\_EqState} & on \texttt{T\_Unit} and \texttt{T\_Tank}: 0=Available-clean 1=In use 2=Dirty 3=Cleaning (CIP) 4=Sterilising (SIP) 5=Maintenance 6=Out of service \\
Mode (\texttt{T\_Mode}) & \texttt{Sts\_Mode} & 0=Out of service (L) 1=Maintenance (MT) 2=Manual (M) 3=Automatic (A). Definitions and allowed changes: Section~\ref{sec:pclib-modes}. \\
Field device health & \texttt{Sts\_NE107} & NAMUR NE107: Good, F (failure), C (function check), S (out of specification), M (maintenance required) \\
\bottomrule
\end{longtable}
\end{small}

\section{Importing into Control Expert and Ignition}

\textbf{Control Expert Classic V16.2.}
\begin{enumerate}[itemsep=1pt]
  \item In the Project Browser, right-click \textbf{Derived Data Types}, choose \textbf{Import}, and select \texttt{export/PCLib\_CE\_V16.2.xdd} (all types in dependency order). If a type already exists, choose whether to keep or replace it.
  \item Run \textbf{Build \textgreater{} Analyze Project}.
  \item Declare one unlocated instance per piece of equipment, for example \texttt{P101 : T\_Motor;} or \texttt{FT101 : T\_AI;}.
  \item For Ignition (or any OPC UA client) to read instances by name, enable the data dictionary: \textbf{Tools \textgreater{} Project Settings \textgreater{} PLC embedded data \textgreater{} Data dictionary}, then rebuild and download.
\end{enumerate}
Totals use \texttt{LREAL}; if the CPU does not support it, regenerate with \texttt{gen\_pclib.py --no-lreal}.

\textbf{Ignition 8.1 Perspective.} \texttt{export/PCLib\_Ignition81.json} holds the matching UDT definitions in a \texttt{PCLib} folder. Nested DDTs (\texttt{Mode}, \texttt{Alm\_*}, \texttt{Run}, \texttt{Drv}) become nested UDT instances, and arrays become folders with one tag per element.
\begin{enumerate}[itemsep=1pt]
  \item In the Designer Tag Browser, select \textbf{Data Types}, choose \textbf{Import Tags}, and select the JSON file.
  \item Create the OPC UA connection to the PLC (M580 embedded server, BMENUA0100, or OPC UA Server Expert; see Section~\ref{sec:protocols}).
  \item Create a UDT instance under \textbf{Tags} (for example \texttt{P101} of type \texttt{PCLib/T\_Motor}) and set its parameters as in Table~\ref{tab:pclib-ign-params}.
  \item Bind Perspective faceplates to the instance. \texttt{Cfg\_}, \texttt{Cmd\_}, \texttt{Sim\_}, and \texttt{Prm\_} tags are writable; everything else is read-only, so the HMI never writes PLC status (ISA-101).
\end{enumerate}


\begin{longtable}{@{}P{3.00cm}P{7.50cm}P{6.43cm}@{}}
\caption{PCLib Ignition UDT Parameters}\label{tab:pclib-ign-params} \\
\toprule
\rowcolor{tblhead}\textbf{Parameter} & \textbf{Meaning} & \textbf{Example} \\
\midrule
\endfirsthead
\multicolumn{3}{@{}l}{\small\textit{Table \thetable{} continued}} \\
\toprule
\rowcolor{tblhead}\textbf{Parameter} & \textbf{Meaning} & \textbf{Example} \\
\midrule
\endhead
\texttt{OPCServer} & Ignition OPC connection name & \texttt{M580\_Line1} \\
\texttt{PathPrefix} & Node-ID prefix used by your OPC server & \texttt{ns=2;s=} (check with the OPC Quick Client) \\
\texttt{BaseTag} & PLC instance name & \texttt{P101} \\
\bottomrule
\end{longtable}

Each tag's OPC path is \texttt{\{PathPrefix\}\{BaseTag\}.\textless{}Member\textgreater{}}. Browse one variable in the OPC Quick Client first and copy its prefix, because the namespace index and path syntax differ between Schneider OPC servers.

\begin{longtable}{@{}P{4.50cm}P{4.00cm}P{8.43cm}@{}}
\caption{PCLib Data Type Mapping}\label{tab:pclib-type-map} \\
\toprule
\rowcolor{tblhead}\textbf{Control Expert} & \textbf{Ignition} & \textbf{Note} \\
\midrule
\endfirsthead
\multicolumn{3}{@{}l}{\small\textit{Table \thetable{} continued}} \\
\toprule
\rowcolor{tblhead}\textbf{Control Expert} & \textbf{Ignition} & \textbf{Note} \\
\midrule
\endhead
BOOL & Boolean &  \\
INT & Int2 &  \\
UINT, WORD, DINT & Int4 &  \\
UDINT & Int8 &  \\
REAL /\allowbreak{} LREAL & Float4 /\allowbreak{} Float8 &  \\
TIME & Int4 & milliseconds (\texttt{engUnit} = ms) \\
DT & DateTime & if your OPC server delivers DT as a number, change these tags to Int8 \\
STRING[n] & String &  \\
\bottomrule
\end{longtable}

\textbf{Alarms:} each \texttt{T\_Alarm} /\allowbreak{} \texttt{T\_AlmLim} UDT has an Ignition alarm on \texttt{Sts\_Active}, so PLC alarms appear in the Perspective Alarm Status Table and alarm journal. The alarm is enabled by \texttt{Cfg\_En}, and its priority comes from \texttt{Cfg\_Priority}. EFC uses Ignition's priority numbering (0=Diagnostic 1=Low 2=Medium 3=High 4=Critical, Table~\ref{tab:alarm-priority}) \newmark.
\note{The member tables below still show the v1.00 numbering (1=Emergency to 4=Low). They are regenerated from \texttt{gen\_pclib.py}, which will be updated to Ignition numbering at the next PCLib revision.}
 Acknowledging in Ignition runs a tag event script that writes \texttt{Cmd\_Ack} to the PLC, so the PLC alarm state (ISA-18.2) stays authoritative. Shelving is done in the PLC through \texttt{Cmd\_Shelve}; Ignition's own shelving only hides the alarm in Ignition.

\section{Type Index}

\begin{small}
\begin{longtable}{@{}P{2.60cm}P{3.70cm}P{3.70cm}P{1.40cm}P{4.69cm}@{}}
\caption{PCLib Type Index}\label{tab:pclib-index} \\
\toprule
\rowcolor{tblhead}\textbf{Type} & \textbf{ISA-88 role} & \textbf{MTP equivalent} & \textbf{Mem.} & \textbf{Depends on} \\
\midrule
\endfirsthead
\multicolumn{5}{@{}l}{\small\textit{Table \thetable{} continued}} \\
\toprule
\rowcolor{tblhead}\textbf{Type} & \textbf{ISA-88 role} & \textbf{MTP equivalent} & \textbf{Mem.} & \textbf{Depends on} \\
\midrule
\endhead
\hyperref[sec:pclib-t-alarm]{\texttt{T\_Alarm}} & Support (alarm object) & Alarm/\allowbreak{}Warning/\allowbreak{}Tolerance limit attributes & 26 & - \\
\hyperref[sec:pclib-t-almlim]{\texttt{T\_AlmLim}} & Support (alarm object) & e.g. VAHLim /\allowbreak{} VAHEn /\allowbreak{} VAHAct & 28 & - \\
\hyperref[sec:pclib-t-mode]{\texttt{T\_Mode}} & Support (mode + ownership) & OpMode (Offline/\allowbreak{}Operator/\allowbreak{}Automatic) + SrcMode & 20 & - \\
\hyperref[sec:pclib-t-runtime]{\texttt{T\_RunTime}} & Support (asset data) & - & 11 & - \\
\hyperref[sec:pclib-t-ai]{\texttt{T\_AI}} & Control module & AnaMon & 35 & T\_Alarm, T\_AlmLim \\
\hyperref[sec:pclib-t-di]{\texttt{T\_DI}} & Control module & BinMon & 19 & T\_Alarm \\
\hyperref[sec:pclib-t-dvalve]{\texttt{T\_DValve}} & Control module & MonBinVlv & 41 & T\_Alarm, T\_Mode, T\_RunTime \\
\hyperref[sec:pclib-t-cvalve]{\texttt{T\_CValve}} & Control module & MonAnaVlv & 42 & T\_Alarm, T\_Mode \\
\hyperref[sec:pclib-t-motor]{\texttt{T\_Motor}} & Control module & MonBinDrv & 52 & T\_Alarm, T\_AlmLim, T\_Mode, T\_RunTime \\
\hyperref[sec:pclib-t-vfd]{\texttt{T\_VFD}} & Control module & MonAnaDrv & 82 & T\_Alarm, T\_Mode, T\_RunTime \\
\hyperref[sec:pclib-t-heater]{\texttt{T\_Heater}} & Control module & MonBinDrv (on/\allowbreak{}off) /\allowbreak{} MonAnaDrv (SSR, SCR) & 64 & T\_Alarm, T\_Mode, T\_RunTime \\
\hyperref[sec:pclib-t-pid]{\texttt{T\_PID}} & Control module & PIDCtrl & 51 & T\_Alarm, T\_AlmLim \\
\hyperref[sec:pclib-t-pumpprot]{\texttt{T\_PumpProt}} & Control module & - & 28 & T\_Alarm, T\_AlmLim \\
\hyperref[sec:pclib-t-pumpds]{\texttt{T\_PumpDS}} & Equipment module & - & 26 & T\_Alarm, T\_Mode \\
\hyperref[sec:pclib-t-agitator]{\texttt{T\_Agitator}} & Equipment module & MonAnaDrv (drive) & 15 & T\_Alarm, T\_AlmLim, T\_Mode, T\_RunTime, T\_VFD \\
\hyperref[sec:pclib-t-compressor]{\texttt{T\_Compressor}} & Equipment module & MonAnaDrv (drive) & 29 & T\_Alarm, T\_AlmLim, T\_Mode, T\_RunTime, T\_VFD \\
\hyperref[sec:pclib-t-conveyor]{\texttt{T\_Conveyor}} & Equipment module & MonBinDrv (drive) & 28 & T\_Alarm, T\_AlmLim, T\_Mode, T\_Motor, T\_RunTime \\
\hyperref[sec:pclib-t-diacomp]{\texttt{T\_DiaComp}} & Equipment module & MonBinDrv (drive) & 89 & T\_Alarm, T\_AlmLim, T\_Mode, T\_Motor \\
\hyperref[sec:pclib-t-tank]{\texttt{T\_Tank}} & Equipment module /\allowbreak{} Unit vessel & - & 41 & T\_Alarm, T\_AlmLim \\
\hyperref[sec:pclib-t-totalizer]{\texttt{T\_Totalizer}} & Control module & - & 41 & T\_Alarm, T\_AlmLim \\
\hyperref[sec:pclib-t-heatex]{\texttt{T\_HeatEx}} & Control module (monitoring) & - & 31 & T\_AlmLim \\
\hyperref[sec:pclib-t-dosepump]{\texttt{T\_DosePump}} & Equipment module & MonAnaDrv (drive) & 31 & T\_Alarm, T\_AlmLim, T\_Mode, T\_Motor, T\_RunTime \\
\hyperref[sec:pclib-t-analyzer]{\texttt{T\_Analyzer}} & Control module & AnaMon & 28 & T\_AI, T\_Alarm, T\_AlmLim \\
\hyperref[sec:pclib-t-gasdet]{\texttt{T\_GasDet}} & Control module & AnaMon & 39 & T\_Alarm, T\_AlmLim \\
\hyperref[sec:pclib-t-intlk]{\texttt{T\_Intlk}} & Support (interlock group) & LockView16 & 26 & T\_Alarm \\
\hyperref[sec:pclib-t-esd]{\texttt{T\_ESD}} & Support (shutdown group) & - & 34 & T\_Alarm \\
\hyperref[sec:pclib-t-phase]{\texttt{T\_Phase}} & Equipment phase (procedural) & Service (state machine) & 42 & T\_Alarm \\
\hyperref[sec:pclib-t-unit]{\texttt{T\_Unit}} & Unit (physical) & - & 26 & T\_Alarm \\
\bottomrule
\end{longtable}
\end{small}

\section{Common Building Blocks}

\subsection{T\_Alarm}\label{sec:pclib-t-alarm}

Discrete alarm, ISA-18.2 /\allowbreak{} IEC 62682 state model, shelving, suppression, out-of-service.

\par\smallskip{\small \textbf{ISA-88 role:} Support (alarm object) \quad \textbf{MTP:} Alarm/\allowbreak{}Warning/\allowbreak{}Tolerance limit attributes \quad \textbf{Standards:} ISA-18.2 /\allowbreak{} IEC 62682, ISA-101.01}\par

\begin{small}
\begin{longtable}{@{}P{3.70cm}P{2.20cm}P{1.70cm}P{8.91cm}@{}}
\caption{T\_Alarm Members}\label{tab:pclib-t-alarm} \\
\toprule
\rowcolor{tblhead}\textbf{Member} & \textbf{Type} & \textbf{Init} & \textbf{Description} \\
\midrule
\endfirsthead
\multicolumn{4}{@{}l}{\small\textit{Table \thetable{} continued}} \\
\toprule
\rowcolor{tblhead}\textbf{Member} & \textbf{Type} & \textbf{Init} & \textbf{Description} \\
\midrule
\endhead
\texttt{Inp\_Cond} & \texttt{BOOL} &  & Raw alarm condition (written by owner logic) \\
\texttt{Inp\_Suppress} & \texttt{BOOL} &  & Suppressed-by-design: state-based suppression by logic (ISA-18.2 state G) \\
\texttt{Cfg\_En} & \texttt{BOOL} & TRUE & Alarm configured (FALSE = not an alarm, e.g. instrument not fitted) \\
\texttt{Cfg\_Latch} & \texttt{BOOL} &  & Latching alarm: active until Cmd\_Reset after RTN (ISA-18.2 state E) \\
\texttt{Cfg\_NoAck} & \texttt{BOOL} &  & No acknowledge required --- use only for alarms rationalised as such \\
\texttt{Cfg\_OnDly} & \texttt{TIME} & T\#0S & On-delay (ISA-18.2 chattering control) \\
\texttt{Cfg\_OffDly} & \texttt{TIME} & T\#0S & Off-delay (ISA-18.2 chattering control) \\
\texttt{Cfg\_Priority} & \texttt{INT} & 3 & Rationalised priority: 1=Emergency 2=High 3=Medium 4=Low \\
\texttt{Cfg\_Class} & \texttt{INT} & 0 & Alarm class (alarm philosophy): 0=General 1=Safety 2=Environmental 3=Quality 4=Asset \\
\texttt{Cfg\_MaxShelf} & \texttt{TIME} & T\#8H & Maximum shelve duration (shelving must be time-limited) \\
\texttt{Cmd\_Ack} & \texttt{BOOL} &  & Operator acknowledge (one-shot) \\
\texttt{Cmd\_Reset} & \texttt{BOOL} &  & Reset latched alarm (one-shot) \\
\texttt{Cmd\_Shelve} & \texttt{BOOL} &  & Operator shelve (one-shot) \\
\texttt{Cmd\_Unshelve} & \texttt{BOOL} &  & Operator unshelve (one-shot) \\
\texttt{Cmd\_RemoveSvc} & \texttt{BOOL} &  & Remove from service for maintenance (one-shot, permit required) \\
\texttt{Cmd\_ReturnSvc} & \texttt{BOOL} &  & Return to service (one-shot) \\
\texttt{Sts\_Active} & \texttt{BOOL} &  & Alarm active (after delays and latch) \\
\texttt{Sts\_Unack} & \texttt{BOOL} &  & Alarm unacknowledged \\
\texttt{Sts\_Shelved} & \texttt{BOOL} &  & Shelved by operator (ISA-18.2 state F) \\
\texttt{Sts\_Suppressed} & \texttt{BOOL} &  & Suppressed by design (ISA-18.2 state G) \\
\texttt{Sts\_OOS} & \texttt{BOOL} &  & Out of service (ISA-18.2 state H) \\
\texttt{Sts\_State} & \texttt{INT} &  & 0=A Normal 1=B Unack 2=C Acked 3=D RTN-Unack 4=E Latched 5=F Shelved 6=G Suppressed-by-design 7=H Out-of-service \\
\texttt{Sts\_ShelfLeft} & \texttt{TIME} &  & Shelve time remaining \\
\texttt{Val\_Count} & \texttt{UDINT} &  & Activation counter (alarm performance KPIs) \\
\texttt{Val\_LastActive} & \texttt{DT} &  & Date/\allowbreak{}time of last activation \\
\texttt{Val\_LastRTN} & \texttt{DT} &  & Date/\allowbreak{}time of last return to normal \\
\bottomrule
\end{longtable}
\end{small}

\subsection{T\_AlmLim}\label{sec:pclib-t-almlim}

Analog limit alarm: limit + deadband + ISA-18.2 /\allowbreak{} IEC 62682 alarm.

\par\smallskip{\small \textbf{ISA-88 role:} Support (alarm object) \quad \textbf{MTP:} e.g. VAHLim /\allowbreak{} VAHEn /\allowbreak{} VAHAct \quad \textbf{Standards:} ISA-18.2 /\allowbreak{} IEC 62682, ISA-101.01}\par

\begin{small}
\begin{longtable}{@{}P{3.70cm}P{2.20cm}P{1.70cm}P{8.91cm}@{}}
\caption{T\_AlmLim Members}\label{tab:pclib-t-almlim} \\
\toprule
\rowcolor{tblhead}\textbf{Member} & \textbf{Type} & \textbf{Init} & \textbf{Description} \\
\midrule
\endfirsthead
\multicolumn{4}{@{}l}{\small\textit{Table \thetable{} continued}} \\
\toprule
\rowcolor{tblhead}\textbf{Member} & \textbf{Type} & \textbf{Init} & \textbf{Description} \\
\midrule
\endhead
\texttt{Cfg\_Lim} & \texttt{REAL} & 0.0 & Alarm limit (EU) --- from rationalisation /\allowbreak{} master alarm database \\
\texttt{Cfg\_DB} & \texttt{REAL} & 0.0 & Deadband for return to normal (EU, ISA-18.2 chattering control) \\
\texttt{Inp\_Cond} & \texttt{BOOL} &  & Raw alarm condition (written by owner logic) \\
\texttt{Inp\_Suppress} & \texttt{BOOL} &  & Suppressed-by-design: state-based suppression by logic (ISA-18.2 state G) \\
\texttt{Cfg\_En} & \texttt{BOOL} &  & Alarm enabled --- set TRUE once Cfg\_Lim is rationalised \\
\texttt{Cfg\_Latch} & \texttt{BOOL} &  & Latching alarm: active until Cmd\_Reset after RTN (ISA-18.2 state E) \\
\texttt{Cfg\_NoAck} & \texttt{BOOL} &  & No acknowledge required --- use only for alarms rationalised as such \\
\texttt{Cfg\_OnDly} & \texttt{TIME} & T\#0S & On-delay (ISA-18.2 chattering control) \\
\texttt{Cfg\_OffDly} & \texttt{TIME} & T\#0S & Off-delay (ISA-18.2 chattering control) \\
\texttt{Cfg\_Priority} & \texttt{INT} & 3 & Rationalised priority: 1=Emergency 2=High 3=Medium 4=Low \\
\texttt{Cfg\_Class} & \texttt{INT} & 0 & Alarm class (alarm philosophy): 0=General 1=Safety 2=Environmental 3=Quality 4=Asset \\
\texttt{Cfg\_MaxShelf} & \texttt{TIME} & T\#8H & Maximum shelve duration (shelving must be time-limited) \\
\texttt{Cmd\_Ack} & \texttt{BOOL} &  & Operator acknowledge (one-shot) \\
\texttt{Cmd\_Reset} & \texttt{BOOL} &  & Reset latched alarm (one-shot) \\
\texttt{Cmd\_Shelve} & \texttt{BOOL} &  & Operator shelve (one-shot) \\
\texttt{Cmd\_Unshelve} & \texttt{BOOL} &  & Operator unshelve (one-shot) \\
\texttt{Cmd\_RemoveSvc} & \texttt{BOOL} &  & Remove from service for maintenance (one-shot, permit required) \\
\texttt{Cmd\_ReturnSvc} & \texttt{BOOL} &  & Return to service (one-shot) \\
\texttt{Sts\_Active} & \texttt{BOOL} &  & Alarm active (after delays and latch) \\
\texttt{Sts\_Unack} & \texttt{BOOL} &  & Alarm unacknowledged \\
\texttt{Sts\_Shelved} & \texttt{BOOL} &  & Shelved by operator (ISA-18.2 state F) \\
\texttt{Sts\_Suppressed} & \texttt{BOOL} &  & Suppressed by design (ISA-18.2 state G) \\
\texttt{Sts\_OOS} & \texttt{BOOL} &  & Out of service (ISA-18.2 state H) \\
\texttt{Sts\_State} & \texttt{INT} &  & 0=A Normal 1=B Unack 2=C Acked 3=D RTN-Unack 4=E Latched 5=F Shelved 6=G Suppressed-by-design 7=H Out-of-service \\
\texttt{Sts\_ShelfLeft} & \texttt{TIME} &  & Shelve time remaining \\
\texttt{Val\_Count} & \texttt{UDINT} &  & Activation counter (alarm performance KPIs) \\
\texttt{Val\_LastActive} & \texttt{DT} &  & Date/\allowbreak{}time of last activation \\
\texttt{Val\_LastRTN} & \texttt{DT} &  & Date/\allowbreak{}time of last return to normal \\
\bottomrule
\end{longtable}
\end{small}

\subsection{T\_Mode}\label{sec:pclib-t-mode}

Control-module mode and ISA-88 ownership (arbitration); maps to MTP OpMode.

\par\smallskip{\small \textbf{ISA-88 role:} Support (mode + ownership) \quad \textbf{MTP:} OpMode (Offline/\allowbreak{}Operator/\allowbreak{}Automatic) + SrcMode \quad \textbf{Standards:} ISA-88 /\allowbreak{} IEC 61512-1, VDI/\allowbreak{}VDE/\allowbreak{}NAMUR 2658-4 (MTP)}\par

\begin{small}
\begin{longtable}{@{}P{3.70cm}P{2.20cm}P{1.70cm}P{8.91cm}@{}}
\caption{T\_Mode Members}\label{tab:pclib-t-mode} \\
\toprule
\rowcolor{tblhead}\textbf{Member} & \textbf{Type} & \textbf{Init} & \textbf{Description} \\
\midrule
\endfirsthead
\multicolumn{4}{@{}l}{\small\textit{Table \thetable{} continued}} \\
\toprule
\rowcolor{tblhead}\textbf{Member} & \textbf{Type} & \textbf{Init} & \textbf{Description} \\
\midrule
\endhead
\texttt{Cmd\_Auto} & \texttt{BOOL} &  & Operator request Automatic (one-shot) \\
\texttt{Cmd\_Man} & \texttt{BOOL} &  & Operator request Manual (one-shot) \\
\texttt{Cmd\_Maint} & \texttt{BOOL} &  & Request Maintenance (one-shot) \\
\texttt{Cmd\_OOS} & \texttt{BOOL} &  & Request Out-of-Service (one-shot) \\
\texttt{Prog\_Auto} & \texttt{BOOL} &  & Owner request Automatic \\
\texttt{Prog\_Man} & \texttt{BOOL} &  & Owner request Manual \\
\texttt{Prog\_Lock} & \texttt{BOOL} &  & Owner locks current mode (operator cannot change) \\
\texttt{Prog\_Acquire} & \texttt{BOOL} &  & ISA-88 arbitration: owner (phase /\allowbreak{} EM) acquires this module \\
\texttt{Prog\_Release} & \texttt{BOOL} &  & ISA-88 arbitration: owner releases this module \\
\texttt{Inp\_ReqOwnerID} & \texttt{DINT} &  & ID of the phase /\allowbreak{} equipment module requesting ownership \\
\texttt{Cmd\_ForceRelease} & \texttt{BOOL} &  & Supervisor forces release of ownership (one-shot) \\
\texttt{Cfg\_PowerUpMode} & \texttt{INT} & 2 & Mode at cold start: 0=OOS 1=Maint 2=Man 3=Auto \\
\texttt{Cfg\_NoMaint} & \texttt{BOOL} &  & Maintenance mode not allowed \\
\texttt{Sts\_Auto} & \texttt{BOOL} &  & In Automatic (commanded by owner) --- MTP Automatic \\
\texttt{Sts\_Man} & \texttt{BOOL} &  & In Manual (commanded by operator) --- MTP Operator \\
\texttt{Sts\_Maint} & \texttt{BOOL} &  & In Maintenance \\
\texttt{Sts\_OOS} & \texttt{BOOL} &  & Out of service --- MTP Offline \\
\texttt{Sts\_Mode} & \texttt{INT} &  & 0=OOS 1=Maint 2=Man 3=Auto  (MTP: 0=Offline 2=Operator 3=Automatic) \\
\texttt{Sts\_Owned} & \texttt{BOOL} &  & Acquired by an owner \\
\texttt{Val\_OwnerID} & \texttt{DINT} &  & ID of current owner, 0 = free \\
\bottomrule
\end{longtable}
\end{small}

\subsection{T\_RunTime}\label{sec:pclib-t-runtime}

Run-hours, starts and maintenance counters.

\par\smallskip{\small \textbf{ISA-88 role:} Support (asset data) \quad \textbf{MTP:} - \quad \textbf{Standards:} NAMUR NE107 (feeds Maintenance required)}\par

\begin{small}
\begin{longtable}{@{}P{3.70cm}P{2.20cm}P{1.70cm}P{8.91cm}@{}}
\caption{T\_RunTime Members}\label{tab:pclib-t-runtime} \\
\toprule
\rowcolor{tblhead}\textbf{Member} & \textbf{Type} & \textbf{Init} & \textbf{Description} \\
\midrule
\endfirsthead
\multicolumn{4}{@{}l}{\small\textit{Table \thetable{} continued}} \\
\toprule
\rowcolor{tblhead}\textbf{Member} & \textbf{Type} & \textbf{Init} & \textbf{Description} \\
\midrule
\endhead
\texttt{Inp\_Running} & \texttt{BOOL} &  & Equipment running (from owner logic) \\
\texttt{Cfg\_MaintHrs} & \texttt{REAL} & 0.0 & Maintenance interval in run hours, 0 = off \\
\texttt{Cmd\_ResetMaint} & \texttt{BOOL} &  & Reset hours-since-maintenance (one-shot) \\
\texttt{Cmd\_ResetAll} & \texttt{BOOL} &  & Reset all counters (one-shot, supervisor) \\
\texttt{Val\_RunSecs} & \texttt{UDINT} &  & Total run time (s) --- master counter \\
\texttt{Val\_RunHrs} & \texttt{REAL} &  & Total run hours (display) \\
\texttt{Val\_Starts} & \texttt{UDINT} &  & Total starts /\allowbreak{} cycles \\
\texttt{Val\_StartsLastHr} & \texttt{INT} &  & Starts in the last 60 min (rolling) \\
\texttt{Val\_LastRun} & \texttt{TIME} &  & Duration of last run \\
\texttt{Val\_HrsSinceMaint} & \texttt{REAL} &  & Run hours since last maintenance reset \\
\texttt{Sts\_MaintDue} & \texttt{BOOL} &  & Maintenance due \\
\bottomrule
\end{longtable}
\end{small}

\section{Core Instruments and Devices}

\subsection{T\_AI}\label{sec:pclib-t-ai}

Analog input: scaling, filtering, NE43 range check, alarms, sim.

\par\smallskip{\small \textbf{ISA-88 role:} Control module \quad \textbf{MTP:} AnaMon \quad \textbf{Standards:} ISA-88 /\allowbreak{} IEC 61512-1, ISA-5.1, ISA-18.2 /\allowbreak{} IEC 62682, ISA-101.01, NAMUR NE107, VDI/\allowbreak{}VDE/\allowbreak{}NAMUR 2658-4 (MTP), NAMUR NE43}\par

\begin{small}
\begin{longtable}{@{}P{3.70cm}P{2.20cm}P{1.70cm}P{8.91cm}@{}}
\caption{T\_AI Members}\label{tab:pclib-t-ai} \\
\toprule
\rowcolor{tblhead}\textbf{Member} & \textbf{Type} & \textbf{Init} & \textbf{Description} \\
\midrule
\endfirsthead
\multicolumn{4}{@{}l}{\small\textit{Table \thetable{} continued}} \\
\toprule
\rowcolor{tblhead}\textbf{Member} & \textbf{Type} & \textbf{Init} & \textbf{Description} \\
\midrule
\endhead
\texttt{Cfg\_Tag} & \texttt{STRING[20]} &  & Equipment /\allowbreak{} instrument tag, e.g. FT-101 \\
\texttt{Cfg\_Desc} & \texttt{STRING[40]} &  & Description shown on HMI faceplate \\
\texttt{Cfg\_EU} & \texttt{STRING[12]} &  & Engineering units text \\
\texttt{Inp\_Raw} & \texttt{INT} &  & Raw counts from analog channel \\
\texttt{Inp\_ChFault} & \texttt{BOOL} &  & Channel /\allowbreak{} module diagnostic fault \\
\texttt{Cfg\_RawMin} & \texttt{INT} & 0 & Raw counts at 0\% (M340/\allowbreak{}M580 4-20mA default 0) \\
\texttt{Cfg\_RawMax} & \texttt{INT} & 10000 & Raw counts at 100\% (M340/\allowbreak{}M580 4-20mA default 10000) \\
\texttt{Cfg\_EUMin} & \texttt{REAL} & 0.0 & PV at 0\% (EU) \\
\texttt{Cfg\_EUMax} & \texttt{REAL} & 100.0 & PV at 100\% (EU) \\
\texttt{Cfg\_FiltTC} & \texttt{REAL} & 0.0 & 1st-order filter time constant (s), 0 = off \\
\texttt{Cfg\_SqrtExtract} & \texttt{BOOL} &  & Square-root extraction (DP flow) \\
\texttt{Cfg\_LoCutoff} & \texttt{REAL} & 0.0 & Low cutoff (\% span): PV forced to EUMin below this, 0 = off \\
\texttt{Cfg\_BadLoPct} & \texttt{REAL} & -2.5 & Bad below this \% span (NE43 3.6mA = -2.5\%) \\
\texttt{Cfg\_BadHiPct} & \texttt{REAL} & 106.25 & Bad above this \% span (NE43 21mA = 106.25\%) \\
\texttt{Cfg\_BadAction} & \texttt{INT} & 0 & On bad PV: 0=Hold last good 1=Substitute 2=EUMin 3=EUMax \\
\texttt{Cfg\_SubstPV} & \texttt{REAL} & 0.0 & Substitute PV used when Cfg\_BadAction = 1 \\
\texttt{Sim\_En} & \texttt{BOOL} &  & Simulation enabled \\
\texttt{Sim\_PV} & \texttt{REAL} & 0.0 & Simulated PV (EU) \\
\texttt{Val\_PV} & \texttt{REAL} &  & Process value (EU) --- use this in logic \\
\texttt{Val\_PVPct} & \texttt{REAL} &  & Process value (\% of span) \\
\texttt{Val\_ROC} & \texttt{REAL} &  & Rate of change (EU/\allowbreak{}s, 1 s window) --- input to Alm\_ROC \\
\texttt{Sts\_Bad} & \texttt{BOOL} &  & PV bad (out of range or channel fault) \\
\texttt{Sts\_URng} & \texttt{BOOL} &  & Under range \\
\texttt{Sts\_ORng} & \texttt{BOOL} &  & Over range \\
\texttt{Sts\_Sim} & \texttt{BOOL} &  & Simulation active \\
\texttt{Sts\_Subst} & \texttt{BOOL} &  & Substitute /\allowbreak{} held value in use \\
\texttt{Sts\_Quality} & \texttt{INT} &  & 0=Good 1=Uncertain (simulation, NE107 C or S) 2=Bad \\
\texttt{Inp\_NE107} & \texttt{INT} & 0 & NE107 status reported by smart device (HART /\allowbreak{} fieldbus) \\
\texttt{Sts\_NE107} & \texttt{INT} &  & NAMUR NE107 status: 0=Good 1=Failure (F) 2=Function check (C) 3=Out of specification (S) 4=Maintenance required (M) \\
\texttt{Alm\_HH} & \texttt{T\_AlmLim} &  & High-high alarm \\
\texttt{Alm\_H} & \texttt{T\_AlmLim} &  & High alarm \\
\texttt{Alm\_L} & \texttt{T\_AlmLim} &  & Low alarm \\
\texttt{Alm\_LL} & \texttt{T\_AlmLim} &  & Low-low alarm \\
\texttt{Alm\_ROC} & \texttt{T\_AlmLim} &  & Rate-of-change alarm (Cfg\_Lim in EU/\allowbreak{}s) \\
\texttt{Alm\_Bad} & \texttt{T\_Alarm} &  & Bad PV alarm \\
\bottomrule
\end{longtable}
\end{small}

\subsection{T\_DI}\label{sec:pclib-t-di}

Digital input: invert, debounce, sim, optional state alarm.

\par\smallskip{\small \textbf{ISA-88 role:} Control module \quad \textbf{MTP:} BinMon \quad \textbf{Standards:} ISA-88 /\allowbreak{} IEC 61512-1, ISA-5.1, ISA-18.2 /\allowbreak{} IEC 62682, ISA-101.01, NAMUR NE107, VDI/\allowbreak{}VDE/\allowbreak{}NAMUR 2658-4 (MTP)}\par

\begin{small}
\begin{longtable}{@{}P{3.70cm}P{2.20cm}P{1.70cm}P{8.91cm}@{}}
\caption{T\_DI Members}\label{tab:pclib-t-di} \\
\toprule
\rowcolor{tblhead}\textbf{Member} & \textbf{Type} & \textbf{Init} & \textbf{Description} \\
\midrule
\endfirsthead
\multicolumn{4}{@{}l}{\small\textit{Table \thetable{} continued}} \\
\toprule
\rowcolor{tblhead}\textbf{Member} & \textbf{Type} & \textbf{Init} & \textbf{Description} \\
\midrule
\endhead
\texttt{Cfg\_Tag} & \texttt{STRING[20]} &  & Equipment /\allowbreak{} instrument tag, e.g. FT-101 \\
\texttt{Cfg\_Desc} & \texttt{STRING[40]} &  & Description shown on HMI faceplate \\
\texttt{Inp\_Raw} & \texttt{BOOL} &  & Raw input from channel \\
\texttt{Inp\_ChFault} & \texttt{BOOL} &  & Channel /\allowbreak{} module diagnostic fault \\
\texttt{Cfg\_Invert} & \texttt{BOOL} &  & Invert raw signal (fail-safe NC contacts) \\
\texttt{Cfg\_OnDly} & \texttt{TIME} & T\#0S & Debounce on-delay \\
\texttt{Cfg\_OffDly} & \texttt{TIME} & T\#0S & Debounce off-delay \\
\texttt{Cfg\_AlmMode} & \texttt{INT} & 0 & 0=No alarm 1=Alarm when Val TRUE 2=Alarm when Val FALSE \\
\texttt{Cfg\_TrueText} & \texttt{STRING[16]} &  & HMI text for TRUE state, e.g. OPEN \\
\texttt{Cfg\_FalseText} & \texttt{STRING[16]} &  & HMI text for FALSE state, e.g. CLOSED \\
\texttt{Sim\_En} & \texttt{BOOL} &  & Simulation enabled \\
\texttt{Sim\_Val} & \texttt{BOOL} &  & Simulated value \\
\texttt{Val} & \texttt{BOOL} &  & Conditioned value --- use this in logic \\
\texttt{Sts\_Bad} & \texttt{BOOL} &  & Channel fault \\
\texttt{Sts\_Sim} & \texttt{BOOL} &  & Simulation active \\
\texttt{Sts\_NE107} & \texttt{INT} &  & NAMUR NE107 status: 0=Good 1=Failure (F) 2=Function check (C) 3=Out of specification (S) 4=Maintenance required (M) \\
\texttt{Val\_Transitions} & \texttt{UDINT} &  & Count of FALSE-\textgreater{}TRUE transitions \\
\texttt{Alm} & \texttt{T\_Alarm} &  & State alarm (see Cfg\_AlmMode) \\
\texttt{Alm\_Bad} & \texttt{T\_Alarm} &  & Channel fault alarm \\
\bottomrule
\end{longtable}
\end{small}

\subsection{T\_DValve}\label{sec:pclib-t-dvalve}

On/\allowbreak{}off valve: single or dual solenoid, limit switches, travel alarms.

\par\smallskip{\small \textbf{ISA-88 role:} Control module \quad \textbf{MTP:} MonBinVlv \quad \textbf{Standards:} ISA-88 /\allowbreak{} IEC 61512-1, ISA-5.1, ISA-18.2 /\allowbreak{} IEC 62682, ISA-101.01, NAMUR NE107, VDI/\allowbreak{}VDE/\allowbreak{}NAMUR 2658-4 (MTP)}\par

\begin{small}
\begin{longtable}{@{}P{3.70cm}P{2.20cm}P{1.70cm}P{8.91cm}@{}}
\caption{T\_DValve Members}\label{tab:pclib-t-dvalve} \\
\toprule
\rowcolor{tblhead}\textbf{Member} & \textbf{Type} & \textbf{Init} & \textbf{Description} \\
\midrule
\endfirsthead
\multicolumn{4}{@{}l}{\small\textit{Table \thetable{} continued}} \\
\toprule
\rowcolor{tblhead}\textbf{Member} & \textbf{Type} & \textbf{Init} & \textbf{Description} \\
\midrule
\endhead
\texttt{Cfg\_Tag} & \texttt{STRING[20]} &  & Equipment /\allowbreak{} instrument tag, e.g. FT-101 \\
\texttt{Cfg\_Desc} & \texttt{STRING[40]} &  & Description shown on HMI faceplate \\
\texttt{Mode} & \texttt{T\_Mode} &  & Mode block \\
\texttt{Inp\_ZSO} & \texttt{BOOL} &  & Open limit switch \\
\texttt{Inp\_ZSC} & \texttt{BOOL} &  & Closed limit switch \\
\texttt{Inp\_Remote} & \texttt{BOOL} &  & Local/\allowbreak{}Remote switch in REMOTE (if fitted) \\
\texttt{Inp\_IntlkOK} & \texttt{BOOL} &  & Interlocks healthy (FALSE drives to interlock position) \\
\texttt{Inp\_PermOK} & \texttt{BOOL} &  & Permissives healthy (TRUE = may leave safe position) \\
\texttt{Inp\_IOFault} & \texttt{BOOL} &  & I/\allowbreak{}O channel or comms fault \\
\texttt{Cfg\_HasZSO} & \texttt{BOOL} & TRUE & Open limit switch fitted \\
\texttt{Cfg\_HasZSC} & \texttt{BOOL} & TRUE & Closed limit switch fitted \\
\texttt{Cfg\_FailOpen} & \texttt{BOOL} &  & Fail position: TRUE = fail open, FALSE = fail closed \\
\texttt{Cfg\_IntlkOpen} & \texttt{BOOL} &  & Interlock drives valve: TRUE = open, FALSE = closed \\
\texttt{Cfg\_DualCoil} & \texttt{BOOL} &  & Dual solenoid (uses Out\_Open and Out\_Close) \\
\texttt{Cfg\_TravelTime} & \texttt{TIME} & T\#10S & Maximum travel time before failure alarm \\
\texttt{Cmd\_Open} & \texttt{BOOL} &  & Operator open (one-shot) \\
\texttt{Cmd\_Close} & \texttt{BOOL} &  & Operator close (one-shot) \\
\texttt{Cmd\_Reset} & \texttt{BOOL} &  & Fault reset (one-shot) \\
\texttt{Prog\_Open} & \texttt{BOOL} &  & Program open request (Auto) \\
\texttt{Prog\_Close} & \texttt{BOOL} &  & Program close request (Auto) \\
\texttt{Out\_Open} & \texttt{BOOL} &  & Solenoid output: open (energise) \\
\texttt{Out\_Close} & \texttt{BOOL} &  & Solenoid output: close (dual-coil only) \\
\texttt{Sim\_En} & \texttt{BOOL} &  & Simulation: ignore field I/\allowbreak{}O, logic generates feedback \\
\texttt{Sts\_Opened} & \texttt{BOOL} &  & Confirmed open \\
\texttt{Sts\_Closed} & \texttt{BOOL} &  & Confirmed closed \\
\texttt{Sts\_Opening} & \texttt{BOOL} &  & Travelling open \\
\texttt{Sts\_Closing} & \texttt{BOOL} &  & Travelling closed \\
\texttt{Sts\_Avail} & \texttt{BOOL} &  & Available to Auto/\allowbreak{}program control \\
\texttt{Sts\_Intlk} & \texttt{BOOL} &  & Interlock not healthy \\
\texttt{Sts\_NotPerm} & \texttt{BOOL} &  & Permissive not healthy \\
\texttt{Sts\_Local} & \texttt{BOOL} &  & In local control \\
\texttt{Sts\_Fault} & \texttt{BOOL} &  & Any fault present \\
\texttt{Sts\_Sim} & \texttt{BOOL} &  & Simulation active \\
\texttt{Sts\_State} & \texttt{INT} &  & 0=Transit/\allowbreak{}unknown 1=Closed 2=Open 3=Opening 4=Closing 5=Fault \\
\texttt{Sts\_NE107} & \texttt{INT} &  & NAMUR NE107 status: 0=Good 1=Failure (F) 2=Function check (C) 3=Out of specification (S) 4=Maintenance required (M) \\
\texttt{Alm\_FailOpen} & \texttt{T\_Alarm} &  & Failed to open within travel time \\
\texttt{Alm\_FailClose} & \texttt{T\_Alarm} &  & Failed to close within travel time \\
\texttt{Alm\_PosConflict} & \texttt{T\_Alarm} &  & Both limit switches made \\
\texttt{Alm\_Uncmd} & \texttt{T\_Alarm} &  & Left commanded position without a command \\
\texttt{Alm\_IntlkTrip} & \texttt{T\_Alarm} &  & Driven to safe position by interlock \\
\texttt{Run} & \texttt{T\_RunTime} &  & Cycle counter (Val\_Starts = operations) \\
\bottomrule
\end{longtable}
\end{small}

\subsection{T\_CValve}\label{sec:pclib-t-cvalve}

Modulating control valve: limits, rate limit, positioner feedback.

\par\smallskip{\small \textbf{ISA-88 role:} Control module \quad \textbf{MTP:} MonAnaVlv \quad \textbf{Standards:} ISA-88 /\allowbreak{} IEC 61512-1, ISA-5.1, ISA-18.2 /\allowbreak{} IEC 62682, ISA-101.01, NAMUR NE107, VDI/\allowbreak{}VDE/\allowbreak{}NAMUR 2658-4 (MTP)}\par

\begin{small}
\begin{longtable}{@{}P{3.70cm}P{2.20cm}P{1.70cm}P{8.91cm}@{}}
\caption{T\_CValve Members}\label{tab:pclib-t-cvalve} \\
\toprule
\rowcolor{tblhead}\textbf{Member} & \textbf{Type} & \textbf{Init} & \textbf{Description} \\
\midrule
\endfirsthead
\multicolumn{4}{@{}l}{\small\textit{Table \thetable{} continued}} \\
\toprule
\rowcolor{tblhead}\textbf{Member} & \textbf{Type} & \textbf{Init} & \textbf{Description} \\
\midrule
\endhead
\texttt{Cfg\_Tag} & \texttt{STRING[20]} &  & Equipment /\allowbreak{} instrument tag, e.g. FT-101 \\
\texttt{Cfg\_Desc} & \texttt{STRING[40]} &  & Description shown on HMI faceplate \\
\texttt{Mode} & \texttt{T\_Mode} &  & Mode block \\
\texttt{Inp\_PosFbk} & \texttt{REAL} &  & Position feedback (\%) \\
\texttt{Inp\_FbkBad} & \texttt{BOOL} &  & Feedback signal bad \\
\texttt{Inp\_ZSC} & \texttt{BOOL} &  & Closed limit switch (optional) \\
\texttt{Inp\_IntlkOK} & \texttt{BOOL} &  & Interlocks healthy (FALSE drives to Cfg\_IntlkPos) \\
\texttt{Inp\_IOFault} & \texttt{BOOL} &  & I/\allowbreak{}O channel or comms fault \\
\texttt{Cfg\_HasFbk} & \texttt{BOOL} &  & Position feedback fitted \\
\texttt{Cfg\_FailPos} & \texttt{REAL} & 0.0 & Fail position (\%): 0 = fail closed, 100 = fail open \\
\texttt{Cfg\_IntlkPos} & \texttt{REAL} & 0.0 & Position driven to on interlock (\%) \\
\texttt{Cfg\_OutMin} & \texttt{REAL} & 0.0 & Minimum output (\%) \\
\texttt{Cfg\_OutMax} & \texttt{REAL} & 100.0 & Maximum output (\%) \\
\texttt{Cfg\_RateUp} & \texttt{REAL} & 0.0 & Opening rate limit (\%/\allowbreak{}s), 0 = none \\
\texttt{Cfg\_RateDn} & \texttt{REAL} & 0.0 & Closing rate limit (\%/\allowbreak{}s), 0 = none \\
\texttt{Cfg\_TightShut} & \texttt{REAL} & 0.0 & Demand below this (\%) forces 0 for tight shut-off \\
\texttt{Cfg\_ReverseAct} & \texttt{BOOL} &  & Reverse-acting output (100\% signal = closed) \\
\texttt{Cfg\_RawMin} & \texttt{INT} & 0 & Analog output raw counts at 0\% \\
\texttt{Cfg\_RawMax} & \texttt{INT} & 10000 & Analog output raw counts at 100\% \\
\texttt{Cfg\_DevLim} & \texttt{REAL} & 5.0 & Position deviation alarm limit (\%) \\
\texttt{Cfg\_DevTime} & \texttt{TIME} & T\#15S & Deviation must persist this long \\
\texttt{Cmd\_ManPos} & \texttt{REAL} & 0.0 & Operator manual demand (\%) \\
\texttt{Cmd\_Reset} & \texttt{BOOL} &  & Fault reset (one-shot) \\
\texttt{Prog\_Pos} & \texttt{REAL} & 0.0 & Program /\allowbreak{} PID demand (\%) --- Auto mode \\
\texttt{Out\_CV} & \texttt{REAL} &  & Output demand (\%) after limits and rate limit \\
\texttt{Out\_CVRaw} & \texttt{INT} &  & Output in raw counts for AO channel \\
\texttt{Sim\_En} & \texttt{BOOL} &  & Simulation: ignore field I/\allowbreak{}O, logic generates feedback \\
\texttt{Val\_Demand} & \texttt{REAL} &  & Active demand before rate limit (\%) \\
\texttt{Val\_Pos} & \texttt{REAL} &  & Position (\%) --- feedback, or demand if no feedback \\
\texttt{Val\_Dev} & \texttt{REAL} &  & Demand minus position (\%) \\
\texttt{Val\_Travel} & \texttt{REAL} &  & Accumulated travel (\% of full stroke, summed) for diagnostics \\
\texttt{Val\_Reversals} & \texttt{UDINT} &  & Direction reversal count \\
\texttt{Sts\_Closed} & \texttt{BOOL} &  & Valve closed \\
\texttt{Sts\_Intlk} & \texttt{BOOL} &  & Interlock active \\
\texttt{Sts\_RateLim} & \texttt{BOOL} &  & Output is being rate limited \\
\texttt{Sts\_Fault} & \texttt{BOOL} &  & Any fault present \\
\texttt{Sts\_Sim} & \texttt{BOOL} &  & Simulation active \\
\texttt{Inp\_NE107} & \texttt{INT} & 0 & NE107 status reported by smart device (HART /\allowbreak{} fieldbus) \\
\texttt{Sts\_NE107} & \texttt{INT} &  & NAMUR NE107 status: 0=Good 1=Failure (F) 2=Function check (C) 3=Out of specification (S) 4=Maintenance required (M) \\
\texttt{Alm\_PosDev} & \texttt{T\_Alarm} &  & Position deviation \\
\texttt{Alm\_FbkBad} & \texttt{T\_Alarm} &  & Feedback bad \\
\texttt{Alm\_IntlkTrip} & \texttt{T\_Alarm} &  & Driven to interlock position \\
\bottomrule
\end{longtable}
\end{small}

\subsection{T\_Motor}\label{sec:pclib-t-motor}

Direct-on-line /\allowbreak{} star-delta /\allowbreak{} soft-start motor.

\par\smallskip{\small \textbf{ISA-88 role:} Control module \quad \textbf{MTP:} MonBinDrv \quad \textbf{Standards:} ISA-88 /\allowbreak{} IEC 61512-1, ISA-5.1, ISA-18.2 /\allowbreak{} IEC 62682, ISA-101.01, NAMUR NE107, VDI/\allowbreak{}VDE/\allowbreak{}NAMUR 2658-4 (MTP)}\par

\begin{small}
\begin{longtable}{@{}P{3.70cm}P{2.20cm}P{1.70cm}P{8.91cm}@{}}
\caption{T\_Motor Members}\label{tab:pclib-t-motor} \\
\toprule
\rowcolor{tblhead}\textbf{Member} & \textbf{Type} & \textbf{Init} & \textbf{Description} \\
\midrule
\endfirsthead
\multicolumn{4}{@{}l}{\small\textit{Table \thetable{} continued}} \\
\toprule
\rowcolor{tblhead}\textbf{Member} & \textbf{Type} & \textbf{Init} & \textbf{Description} \\
\midrule
\endhead
\texttt{Cfg\_Tag} & \texttt{STRING[20]} &  & Equipment /\allowbreak{} instrument tag, e.g. FT-101 \\
\texttt{Cfg\_Desc} & \texttt{STRING[40]} &  & Description shown on HMI faceplate \\
\texttt{Mode} & \texttt{T\_Mode} &  & Mode block \\
\texttt{Inp\_RunFbk} & \texttt{BOOL} &  & Run feedback (contactor aux /\allowbreak{} drive running) \\
\texttt{Inp\_Remote} & \texttt{BOOL} &  & Local/\allowbreak{}Remote switch in REMOTE \\
\texttt{Inp\_AvailOK} & \texttt{BOOL} &  & MCC available: breaker closed, control power OK \\
\texttt{Inp\_EStopOK} & \texttt{BOOL} &  & Field E-stop circuit healthy \\
\texttt{Inp\_IntlkOK} & \texttt{BOOL} &  & Interlocks healthy (TRUE = may run; FALSE trips) \\
\texttt{Inp\_PermOK} & \texttt{BOOL} &  & Permissives healthy (TRUE = may start) \\
\texttt{Inp\_IOFault} & \texttt{BOOL} &  & I/\allowbreak{}O channel or comms fault \\
\texttt{Inp\_OL} & \texttt{BOOL} &  & Overload /\allowbreak{} motor protection relay tripped \\
\texttt{Inp\_Current} & \texttt{REAL} &  & Motor current (A), optional \\
\texttt{Cfg\_FbkTime} & \texttt{TIME} & T\#5S & Time allowed for run feedback after start/\allowbreak{}stop \\
\texttt{Cfg\_RestartDly} & \texttt{TIME} & T\#0S & Minimum stopped time before restart \\
\texttt{Cfg\_NoFbk} & \texttt{BOOL} &  & No run feedback wired (feedback follows output) \\
\texttt{Cfg\_AutoRestart} & \texttt{BOOL} &  & Auto mode restarts after power loss /\allowbreak{} interlock clears \\
\texttt{Cfg\_RatedAmps} & \texttt{REAL} & 0.0 & Motor nameplate full-load current (A) \\
\texttt{Cfg\_MaxStartsHr} & \texttt{INT} & 0 & Maximum starts per hour, 0 = no limit \\
\texttt{Cmd\_Start} & \texttt{BOOL} &  & Operator start (one-shot) \\
\texttt{Cmd\_Stop} & \texttt{BOOL} &  & Operator stop (one-shot) \\
\texttt{Cmd\_Reset} & \texttt{BOOL} &  & Fault /\allowbreak{} trip reset (one-shot) \\
\texttt{Cmd\_Jog} & \texttt{BOOL} &  & Maintenance jog (maintained, Maint mode only) \\
\texttt{Prog\_Start} & \texttt{BOOL} &  & Program /\allowbreak{} sequence start request (Auto mode) \\
\texttt{Prog\_Stop} & \texttt{BOOL} &  & Program /\allowbreak{} sequence stop request (Auto mode) \\
\texttt{Out\_Run} & \texttt{BOOL} &  & Run output to contactor /\allowbreak{} starter \\
\texttt{Sim\_En} & \texttt{BOOL} &  & Simulation: ignore field I/\allowbreak{}O, logic generates feedback \\
\texttt{Sts\_Running} & \texttt{BOOL} &  & Running (feedback confirmed) \\
\texttt{Sts\_Stopped} & \texttt{BOOL} &  & Stopped (feedback confirmed) \\
\texttt{Sts\_Starting} & \texttt{BOOL} &  & Start commanded, waiting for feedback \\
\texttt{Sts\_Stopping} & \texttt{BOOL} &  & Stop commanded, waiting for feedback \\
\texttt{Sts\_Avail} & \texttt{BOOL} &  & Available to Auto/\allowbreak{}program control \\
\texttt{Sts\_Ready} & \texttt{BOOL} &  & Ready to start now (all permissives, no trip) \\
\texttt{Sts\_Intlk} & \texttt{BOOL} &  & Interlock not healthy \\
\texttt{Sts\_NotPerm} & \texttt{BOOL} &  & Permissive not healthy \\
\texttt{Sts\_Local} & \texttt{BOOL} &  & In local control \\
\texttt{Sts\_Tripped} & \texttt{BOOL} &  & Tripped --- reset required \\
\texttt{Sts\_Fault} & \texttt{BOOL} &  & Any fault present \\
\texttt{Sts\_Sim} & \texttt{BOOL} &  & Simulation active \\
\texttt{Sts\_State} & \texttt{INT} &  & 0=Stopped 1=Starting 2=Running 3=Stopping 4=Tripped 5=Not available \\
\texttt{Sts\_RestartLeft} & \texttt{TIME} &  & Restart-delay time remaining \\
\texttt{Sts\_NE107} & \texttt{INT} &  & NAMUR NE107 status: 0=Good 1=Failure (F) 2=Function check (C) 3=Out of specification (S) 4=Maintenance required (M) \\
\texttt{Val\_Current} & \texttt{REAL} &  & Motor current (A) \\
\texttt{Val\_LoadPct} & \texttt{REAL} &  & Load (\% of Cfg\_RatedAmps) \\
\texttt{Alm\_FailStart} & \texttt{T\_Alarm} &  & No run feedback within Cfg\_FbkTime \\
\texttt{Alm\_FailStop} & \texttt{T\_Alarm} &  & Run feedback still present after stop \\
\texttt{Alm\_Uncmd} & \texttt{T\_Alarm} &  & Running without a command \\
\texttt{Alm\_EStop} & \texttt{T\_Alarm} &  & E-stop activated \\
\texttt{Alm\_IntlkTrip} & \texttt{T\_Alarm} &  & Tripped by interlock \\
\texttt{Alm\_ExcessStarts} & \texttt{T\_Alarm} &  & Start blocked: Cfg\_MaxStartsHr exceeded \\
\texttt{Alm\_OL} & \texttt{T\_Alarm} &  & Overload trip \\
\texttt{Alm\_HiCurrent} & \texttt{T\_AlmLim} &  & High current (A) \\
\texttt{Run} & \texttt{T\_RunTime} &  & Run hours /\allowbreak{} starts \\
\bottomrule
\end{longtable}
\end{small}

\subsection{T\_VFD}\label{sec:pclib-t-vfd}

Variable speed drive (fieldbus or hardwired).

\par\smallskip{\small \textbf{ISA-88 role:} Control module \quad \textbf{MTP:} MonAnaDrv \quad \textbf{Standards:} ISA-88 /\allowbreak{} IEC 61512-1, ISA-5.1, ISA-18.2 /\allowbreak{} IEC 62682, ISA-101.01, NAMUR NE107, VDI/\allowbreak{}VDE/\allowbreak{}NAMUR 2658-4 (MTP)}\par

\begin{small}
\begin{longtable}{@{}P{3.70cm}P{2.20cm}P{1.70cm}P{8.91cm}@{}}
\caption{T\_VFD Members}\label{tab:pclib-t-vfd} \\
\toprule
\rowcolor{tblhead}\textbf{Member} & \textbf{Type} & \textbf{Init} & \textbf{Description} \\
\midrule
\endfirsthead
\multicolumn{4}{@{}l}{\small\textit{Table \thetable{} continued}} \\
\toprule
\rowcolor{tblhead}\textbf{Member} & \textbf{Type} & \textbf{Init} & \textbf{Description} \\
\midrule
\endhead
\texttt{Cfg\_Tag} & \texttt{STRING[20]} &  & Equipment /\allowbreak{} instrument tag, e.g. FT-101 \\
\texttt{Cfg\_Desc} & \texttt{STRING[40]} &  & Description shown on HMI faceplate \\
\texttt{Mode} & \texttt{T\_Mode} &  & Mode block \\
\texttt{Inp\_RunFbk} & \texttt{BOOL} &  & Run feedback (contactor aux /\allowbreak{} drive running) \\
\texttt{Inp\_Remote} & \texttt{BOOL} &  & Local/\allowbreak{}Remote switch in REMOTE \\
\texttt{Inp\_AvailOK} & \texttt{BOOL} &  & MCC available: breaker closed, control power OK \\
\texttt{Inp\_EStopOK} & \texttt{BOOL} &  & Field E-stop circuit healthy \\
\texttt{Inp\_IntlkOK} & \texttt{BOOL} &  & Interlocks healthy (TRUE = may run; FALSE trips) \\
\texttt{Inp\_PermOK} & \texttt{BOOL} &  & Permissives healthy (TRUE = may start) \\
\texttt{Inp\_IOFault} & \texttt{BOOL} &  & I/\allowbreak{}O channel or comms fault \\
\texttt{Inp\_DrvReady} & \texttt{BOOL} &  & Drive ready \\
\texttt{Inp\_DrvFault} & \texttt{BOOL} &  & Drive faulted \\
\texttt{Inp\_DrvWarn} & \texttt{BOOL} &  & Drive warning /\allowbreak{} alarm \\
\texttt{Inp\_FaultCode} & \texttt{UINT} &  & Drive fault code \\
\texttt{Inp\_CommOK} & \texttt{BOOL} & TRUE & Drive communications healthy \\
\texttt{Inp\_NE107} & \texttt{INT} & 0 & NE107 status reported by smart device (HART /\allowbreak{} fieldbus) \\
\texttt{Inp\_SpdFbk} & \texttt{REAL} &  & Speed feedback (Cfg\_SpdEU) \\
\texttt{Inp\_Current} & \texttt{REAL} &  & Motor current (A) \\
\texttt{Inp\_Torque} & \texttt{REAL} &  & Motor torque (\%) \\
\texttt{Inp\_Power} & \texttt{REAL} &  & Motor power (kW) \\
\texttt{Cfg\_FbkTime} & \texttt{TIME} & T\#5S & Time allowed for run feedback after start/\allowbreak{}stop \\
\texttt{Cfg\_RestartDly} & \texttt{TIME} & T\#0S & Minimum stopped time before restart \\
\texttt{Cfg\_NoFbk} & \texttt{BOOL} &  & No run feedback wired (feedback follows output) \\
\texttt{Cfg\_AutoRestart} & \texttt{BOOL} &  & Auto mode restarts after power loss /\allowbreak{} interlock clears \\
\texttt{Cfg\_RatedAmps} & \texttt{REAL} & 0.0 & Motor nameplate full-load current (A) \\
\texttt{Cfg\_MaxStartsHr} & \texttt{INT} & 0 & Maximum starts per hour, 0 = no limit \\
\texttt{Cfg\_SpdEU} & \texttt{STRING[8]} &  & Speed units text (\%, Hz, rpm) \\
\texttt{Cfg\_SpdMin} & \texttt{REAL} & 0.0 & Minimum speed reference \\
\texttt{Cfg\_SpdMax} & \texttt{REAL} & 100.0 & Maximum speed reference \\
\texttt{Cfg\_SpdDevLim} & \texttt{REAL} & 5.0 & Speed deviation alarm limit \\
\texttt{Cfg\_SpdDevTime} & \texttt{TIME} & T\#20S & Speed deviation delay \\
\texttt{Cfg\_Reversible} & \texttt{BOOL} &  & Reverse direction allowed \\
\texttt{Cmd\_Start} & \texttt{BOOL} &  & Operator start (one-shot) \\
\texttt{Cmd\_Stop} & \texttt{BOOL} &  & Operator stop (one-shot) \\
\texttt{Cmd\_Reset} & \texttt{BOOL} &  & Fault /\allowbreak{} trip reset (one-shot) \\
\texttt{Cmd\_Jog} & \texttt{BOOL} &  & Maintenance jog (maintained, Maint mode only) \\
\texttt{Prog\_Start} & \texttt{BOOL} &  & Program /\allowbreak{} sequence start request (Auto mode) \\
\texttt{Prog\_Stop} & \texttt{BOOL} &  & Program /\allowbreak{} sequence stop request (Auto mode) \\
\texttt{Cmd\_Spd} & \texttt{REAL} & 0.0 & Operator manual speed reference \\
\texttt{Cmd\_Rev} & \texttt{BOOL} &  & Operator direction: TRUE = reverse \\
\texttt{Prog\_Spd} & \texttt{REAL} & 0.0 & Program /\allowbreak{} PID speed reference (Auto) \\
\texttt{Prog\_Rev} & \texttt{BOOL} &  & Program direction: TRUE = reverse \\
\texttt{Out\_Run} & \texttt{BOOL} &  & Run command to drive \\
\texttt{Out\_Rev} & \texttt{BOOL} &  & Reverse command to drive \\
\texttt{Out\_SpdRef} & \texttt{REAL} &  & Speed reference to drive (limited) \\
\texttt{Out\_FltReset} & \texttt{BOOL} &  & Fault reset pulse to drive \\
\texttt{Sim\_En} & \texttt{BOOL} &  & Simulation: ignore field I/\allowbreak{}O, logic generates feedback \\
\texttt{Sts\_Running} & \texttt{BOOL} &  & Running (feedback confirmed) \\
\texttt{Sts\_Stopped} & \texttt{BOOL} &  & Stopped (feedback confirmed) \\
\texttt{Sts\_Starting} & \texttt{BOOL} &  & Start commanded, waiting for feedback \\
\texttt{Sts\_Stopping} & \texttt{BOOL} &  & Stop commanded, waiting for feedback \\
\texttt{Sts\_Avail} & \texttt{BOOL} &  & Available to Auto/\allowbreak{}program control \\
\texttt{Sts\_Ready} & \texttt{BOOL} &  & Ready to start now (all permissives, no trip) \\
\texttt{Sts\_Intlk} & \texttt{BOOL} &  & Interlock not healthy \\
\texttt{Sts\_NotPerm} & \texttt{BOOL} &  & Permissive not healthy \\
\texttt{Sts\_Local} & \texttt{BOOL} &  & In local control \\
\texttt{Sts\_Tripped} & \texttt{BOOL} &  & Tripped --- reset required \\
\texttt{Sts\_Fault} & \texttt{BOOL} &  & Any fault present \\
\texttt{Sts\_Sim} & \texttt{BOOL} &  & Simulation active \\
\texttt{Sts\_State} & \texttt{INT} &  & 0=Stopped 1=Starting 2=Running 3=Stopping 4=Tripped 5=Not available \\
\texttt{Sts\_RestartLeft} & \texttt{TIME} &  & Restart-delay time remaining \\
\texttt{Sts\_NE107} & \texttt{INT} &  & NAMUR NE107 status: 0=Good 1=Failure (F) 2=Function check (C) 3=Out of specification (S) 4=Maintenance required (M) \\
\texttt{Sts\_AtSpd} & \texttt{BOOL} &  & At speed (feedback within deviation limit) \\
\texttt{Sts\_Rev} & \texttt{BOOL} &  & Running in reverse \\
\texttt{Sts\_CommFail} & \texttt{BOOL} &  & Communication failure \\
\texttt{Val\_SpdRef} & \texttt{REAL} &  & Active speed reference \\
\texttt{Val\_Spd} & \texttt{REAL} &  & Speed \\
\texttt{Val\_Current} & \texttt{REAL} &  & Motor current (A) \\
\texttt{Val\_Torque} & \texttt{REAL} &  & Torque (\%) \\
\texttt{Val\_Power} & \texttt{REAL} &  & Power (kW) \\
\texttt{Val\_Energy} & \texttt{LREAL} &  & Energy (kWh), integrated from power \\
\texttt{Alm\_FailStart} & \texttt{T\_Alarm} &  & No run feedback within Cfg\_FbkTime \\
\texttt{Alm\_FailStop} & \texttt{T\_Alarm} &  & Run feedback still present after stop \\
\texttt{Alm\_Uncmd} & \texttt{T\_Alarm} &  & Running without a command \\
\texttt{Alm\_EStop} & \texttt{T\_Alarm} &  & E-stop activated \\
\texttt{Alm\_IntlkTrip} & \texttt{T\_Alarm} &  & Tripped by interlock \\
\texttt{Alm\_ExcessStarts} & \texttt{T\_Alarm} &  & Start blocked: Cfg\_MaxStartsHr exceeded \\
\texttt{Alm\_DrvFault} & \texttt{T\_Alarm} &  & Drive fault \\
\texttt{Alm\_DrvWarn} & \texttt{T\_Alarm} &  & Drive warning \\
\texttt{Alm\_CommFail} & \texttt{T\_Alarm} &  & Drive comms failure \\
\texttt{Alm\_SpdDev} & \texttt{T\_Alarm} &  & Speed deviation \\
\texttt{Run} & \texttt{T\_RunTime} &  & Run hours /\allowbreak{} starts \\
\bottomrule
\end{longtable}
\end{small}

\subsection{T\_Heater}\label{sec:pclib-t-heater}

Electric heater: on/\allowbreak{}off contactor, time-proportioned SSR, or SCR power controller.

\par\smallskip{\small \textbf{ISA-88 role:} Control module \quad \textbf{MTP:} MonBinDrv (on/\allowbreak{}off) /\allowbreak{} MonAnaDrv (SSR, SCR) \quad \textbf{Standards:} ISA-88 /\allowbreak{} IEC 61512-1, ISA-5.1, ISA-18.2 /\allowbreak{} IEC 62682, ISA-101.01, NAMUR NE107, VDI/\allowbreak{}VDE/\allowbreak{}NAMUR 2658-4 (MTP), IEC 60947-4-3 (semiconductor controllers, non-motor loads), NFPA 86 (excess-temperature limit)}\par

\begin{small}
\begin{longtable}{@{}P{3.70cm}P{2.20cm}P{1.70cm}P{8.91cm}@{}}
\caption{T\_Heater Members}\label{tab:pclib-t-heater} \\
\toprule
\rowcolor{tblhead}\textbf{Member} & \textbf{Type} & \textbf{Init} & \textbf{Description} \\
\midrule
\endfirsthead
\multicolumn{4}{@{}l}{\small\textit{Table \thetable{} continued}} \\
\toprule
\rowcolor{tblhead}\textbf{Member} & \textbf{Type} & \textbf{Init} & \textbf{Description} \\
\midrule
\endhead
\texttt{Cfg\_Tag} & \texttt{STRING[20]} &  & Equipment /\allowbreak{} instrument tag, e.g. FT-101 \\
\texttt{Cfg\_Desc} & \texttt{STRING[40]} &  & Description shown on HMI faceplate \\
\texttt{Mode} & \texttt{T\_Mode} &  & Mode block \\
\texttt{Inp\_ContactorFbk} & \texttt{BOOL} &  & Main /\allowbreak{} isolating contactor closed (aux contact) \\
\texttt{Inp\_Remote} & \texttt{BOOL} &  & Local/\allowbreak{}Remote switch in REMOTE \\
\texttt{Inp\_AvailOK} & \texttt{BOOL} &  & Supply available: breaker closed, control power OK \\
\texttt{Inp\_HiLimitOK} & \texttt{BOOL} &  & Independent excess-temperature limit healthy (hardwired, latching; PLC monitors only) \\
\texttt{Inp\_IntlkOK} & \texttt{BOOL} &  & Interlocks healthy (TRUE = may heat; FALSE trips) \\
\texttt{Inp\_PermOK} & \texttt{BOOL} &  & Permissives healthy (TRUE = may switch on) \\
\texttt{Inp\_IOFault} & \texttt{BOOL} &  & I/\allowbreak{}O channel or comms fault \\
\texttt{Inp\_Current} & \texttt{REAL} &  & Heater current from CT or SCR controller (A), optional \\
\texttt{Inp\_ScrAlarm} & \texttt{BOOL} &  & SCR controller alarm relay (fuse /\allowbreak{} partial load failure /\allowbreak{} overtemp) \\
\texttt{Cfg\_CtrlType} & \texttt{INT} & 0 & 0 = on/\allowbreak{}off (contactor), 1 = time-proportioned SSR on a PLC output, 2 = SCR power controller (analog reference) \\
\texttt{Cfg\_CycleTime} & \texttt{TIME} & T\#4S & Time-proportioning cycle (type 1): output is ON for Val\_Duty \% of each cycle \\
\texttt{Cfg\_MinPulse} & \texttt{TIME} & T\#100MS & Shortest ON or OFF pulse (type 1); shorter pulses are skipped \\
\texttt{Cfg\_ScrFiring} & \texttt{INT} & 0 & SCR firing mode set in the controller (type 2, documentation): 0 = burst /\allowbreak{} zero-cross, 1 = single-cycle, 2 = phase angle \\
\texttt{Cfg\_PwrMin} & \texttt{REAL} & 0.0 & Minimum power demand when on (\%) \\
\texttt{Cfg\_PwrMax} & \texttt{REAL} & 100.0 & Maximum power demand (\%) \\
\texttt{Cfg\_PwrRate} & \texttt{REAL} & 0.0 & Power ramp limit (\%/\allowbreak{}s), 0 = none (soft start for cold elements) \\
\texttt{Cfg\_RatedKW} & \texttt{REAL} & 0.0 & Heater rated power at 100 \% (kW) \\
\texttt{Cfg\_RatedAmps} & \texttt{REAL} & 0.0 & Heater current at 100 \% (A); 0 = no current monitoring \\
\texttt{Cfg\_LoadFailPct} & \texttt{REAL} & 80.0 & Partial load failure: current below this \% of expected while heating \\
\texttt{Cfg\_LeakAmps} & \texttt{REAL} & 1.0 & Current above this while the output is off = SSR /\allowbreak{} SCR shorted (A) \\
\texttt{Cfg\_FbkTime} & \texttt{TIME} & T\#3S & Time allowed for contactor feedback and current to settle \\
\texttt{Cmd\_On} & \texttt{BOOL} &  & Operator on (one-shot) \\
\texttt{Cmd\_Off} & \texttt{BOOL} &  & Operator off (one-shot) \\
\texttt{Cmd\_Reset} & \texttt{BOOL} &  & Fault /\allowbreak{} trip reset (one-shot) \\
\texttt{Cmd\_Pwr} & \texttt{REAL} &  & Operator power demand (\%), Manual mode, types 1 and 2 \\
\texttt{Prog\_On} & \texttt{BOOL} &  & Program /\allowbreak{} sequence on request (Auto mode) \\
\texttt{Prog\_Off} & \texttt{BOOL} &  & Program /\allowbreak{} sequence off request (Auto mode) \\
\texttt{Prog\_Pwr} & \texttt{REAL} &  & Program power demand (\%), Auto mode, e.g. a T\_PID Out\_CV \\
\texttt{Out\_Contactor} & \texttt{BOOL} &  & Main /\allowbreak{} isolating contactor (opened on trip, high limit or shorted SSR) \\
\texttt{Out\_SSR} & \texttt{BOOL} &  & Time-proportioned output to the SSR (type 1) \\
\texttt{Out\_PwrRef} & \texttt{REAL} &  & Power reference to the SCR controller, 0-100 \% (type 2) \\
\texttt{Sim\_En} & \texttt{BOOL} &  & Simulation: ignore field I/\allowbreak{}O, logic generates feedback \\
\texttt{Sts\_On} & \texttt{BOOL} &  & On: contactor closed and enabled to heat \\
\texttt{Sts\_Off} & \texttt{BOOL} &  & Off \\
\texttt{Sts\_Heating} & \texttt{BOOL} &  & Output conducting now (on/\allowbreak{}off: on; SSR: inside the ON part of the cycle; SCR: reference \textgreater{} 0) \\
\texttt{Sts\_Avail} & \texttt{BOOL} &  & Available to Auto/\allowbreak{}program control \\
\texttt{Sts\_Intlk} & \texttt{BOOL} &  & Interlock not healthy \\
\texttt{Sts\_NotPerm} & \texttt{BOOL} &  & Permissive not healthy \\
\texttt{Sts\_Local} & \texttt{BOOL} &  & In local control \\
\texttt{Sts\_HiLimit} & \texttt{BOOL} &  & Excess-temperature limit tripped (reset at the limit device) \\
\texttt{Sts\_LoadFail} & \texttt{BOOL} &  & Partial load failure (element, fuse or phase lost) \\
\texttt{Sts\_Shorted} & \texttt{BOOL} &  & Current while output off: SSR /\allowbreak{} SCR failed shorted, contactor opened \\
\texttt{Sts\_Tripped} & \texttt{BOOL} &  & Tripped --- reset required \\
\texttt{Sts\_Fault} & \texttt{BOOL} &  & Any fault present \\
\texttt{Sts\_Sim} & \texttt{BOOL} &  & Simulation active \\
\texttt{Sts\_State} & \texttt{INT} &  & 0=Off 1=Switching on 2=On 3=Switching off 4=Tripped 5=Not available \\
\texttt{Sts\_NE107} & \texttt{INT} &  & NAMUR NE107 status: 0=Good 1=Failure (F) 2=Function check (C) 3=Out of specification (S) 4=Maintenance required (M) \\
\texttt{Val\_PwrDemand} & \texttt{REAL} &  & Power demand after limits and ramp (\%) \\
\texttt{Val\_Duty} & \texttt{REAL} &  & Applied duty /\allowbreak{} power (\%): on/\allowbreak{}off 0 or 100, SSR ON-time fraction, SCR reference \\
\texttt{Val\_Power} & \texttt{REAL} &  & Estimated power (kW) = Val\_Duty x Cfg\_RatedKW, or from current \\
\texttt{Val\_Current} & \texttt{REAL} &  & Heater current (A) \\
\texttt{Val\_LoadPct} & \texttt{REAL} &  & Measured current as \% of expected at the present duty \\
\texttt{Val\_Energy} & \texttt{LREAL} &  & Energy total (kWh) \\
\texttt{Alm\_HiLimit} & \texttt{T\_Alarm} &  & Excess-temperature limit tripped \\
\texttt{Alm\_Shorted} & \texttt{T\_Alarm} &  & SSR /\allowbreak{} SCR shorted --- current with output off \\
\texttt{Alm\_LoadFail} & \texttt{T\_Alarm} &  & Partial load failure \\
\texttt{Alm\_FailOn} & \texttt{T\_Alarm} &  & No contactor feedback /\allowbreak{} current within Cfg\_FbkTime \\
\texttt{Alm\_FailOff} & \texttt{T\_Alarm} &  & Contactor feedback still present after off \\
\texttt{Alm\_ScrFault} & \texttt{T\_Alarm} &  & SCR controller alarm \\
\texttt{Alm\_IntlkTrip} & \texttt{T\_Alarm} &  & Tripped by interlock \\
\texttt{Run} & \texttt{T\_RunTime} &  & On hours /\allowbreak{} switch-on count \\
\bottomrule
\end{longtable}
\end{small}

\subsection{T\_PID}\label{sec:pclib-t-pid}

PID loop wrapper (maps onto PIDFF EFB), Man/\allowbreak{}Auto/\allowbreak{}Cas, tracking.

\par\smallskip{\small \textbf{ISA-88 role:} Control module \quad \textbf{MTP:} PIDCtrl \quad \textbf{Standards:} ISA-88 /\allowbreak{} IEC 61512-1, ISA-5.1, ISA-18.2 /\allowbreak{} IEC 62682, ISA-101.01, VDI/\allowbreak{}VDE/\allowbreak{}NAMUR 2658-4 (MTP)}\par

\begin{small}
\begin{longtable}{@{}P{3.70cm}P{2.20cm}P{1.70cm}P{8.91cm}@{}}
\caption{T\_PID Members}\label{tab:pclib-t-pid} \\
\toprule
\rowcolor{tblhead}\textbf{Member} & \textbf{Type} & \textbf{Init} & \textbf{Description} \\
\midrule
\endfirsthead
\multicolumn{4}{@{}l}{\small\textit{Table \thetable{} continued}} \\
\toprule
\rowcolor{tblhead}\textbf{Member} & \textbf{Type} & \textbf{Init} & \textbf{Description} \\
\midrule
\endhead
\texttt{Cfg\_Tag} & \texttt{STRING[20]} &  & Equipment /\allowbreak{} instrument tag, e.g. FT-101 \\
\texttt{Cfg\_Desc} & \texttt{STRING[40]} &  & Description shown on HMI faceplate \\
\texttt{Cfg\_PVEU} & \texttt{STRING[12]} &  & PV engineering units text \\
\texttt{Inp\_PV} & \texttt{REAL} &  & Process value (EU) \\
\texttt{Inp\_PVBad} & \texttt{BOOL} &  & PV bad quality \\
\texttt{Inp\_CasSP} & \texttt{REAL} &  & Cascade (remote) setpoint from master loop \\
\texttt{Inp\_FF} & \texttt{REAL} &  & Feed-forward input (\% CV) \\
\texttt{Inp\_TrackEn} & \texttt{BOOL} &  & Output tracking enable \\
\texttt{Inp\_TrackVal} & \texttt{REAL} &  & Output tracking value (\%) \\
\texttt{Inp\_IntlkOK} & \texttt{BOOL} &  & Interlocks healthy (FALSE forces Cfg\_IntlkCV) \\
\texttt{Cmd\_Man} & \texttt{BOOL} &  & Operator request Manual (one-shot) \\
\texttt{Cmd\_Auto} & \texttt{BOOL} &  & Operator request Auto (one-shot) \\
\texttt{Cmd\_Cas} & \texttt{BOOL} &  & Operator request Cascade (one-shot) \\
\texttt{Cmd\_SP} & \texttt{REAL} & 0.0 & Operator setpoint (EU) \\
\texttt{Cmd\_CV} & \texttt{REAL} & 0.0 & Operator manual output (\%) \\
\texttt{Prog\_SP} & \texttt{REAL} & 0.0 & Program setpoint (EU), used when Cfg\_ProgSP \\
\texttt{Cfg\_ProgSP} & \texttt{BOOL} &  & Setpoint owned by program instead of operator \\
\texttt{Cfg\_Kp} & \texttt{REAL} & 1.0 & Proportional gain \\
\texttt{Cfg\_Ti} & \texttt{TIME} & T\#10S & Integral time, T\#0S = no integral \\
\texttt{Cfg\_Td} & \texttt{TIME} & T\#0S & Derivative time \\
\texttt{Cfg\_DirectAct} & \texttt{BOOL} &  & TRUE = direct acting (CV rises with PV); FALSE = reverse \\
\texttt{Cfg\_PVMin} & \texttt{REAL} & 0.0 & PV range low (EU) \\
\texttt{Cfg\_PVMax} & \texttt{REAL} & 100.0 & PV range high (EU) \\
\texttt{Cfg\_SPMin} & \texttt{REAL} & 0.0 & Setpoint low limit (EU) \\
\texttt{Cfg\_SPMax} & \texttt{REAL} & 100.0 & Setpoint high limit (EU) \\
\texttt{Cfg\_CVMin} & \texttt{REAL} & 0.0 & Output low limit (\%) \\
\texttt{Cfg\_CVMax} & \texttt{REAL} & 100.0 & Output high limit (\%) \\
\texttt{Cfg\_SPRate} & \texttt{REAL} & 0.0 & Setpoint ramp rate (EU/\allowbreak{}s), 0 = step \\
\texttt{Cfg\_CVRate} & \texttt{REAL} & 0.0 & Output rate limit (\%/\allowbreak{}s), 0 = none \\
\texttt{Cfg\_DB} & \texttt{REAL} & 0.0 & Error deadband (EU) \\
\texttt{Cfg\_FFGain} & \texttt{REAL} & 0.0 & Feed-forward gain \\
\texttt{Cfg\_PVTrack} & \texttt{BOOL} &  & Setpoint tracks PV while in Manual (bumpless) \\
\texttt{Cfg\_IntlkCV} & \texttt{REAL} & 0.0 & Output forced on interlock (\%) \\
\texttt{Cfg\_BadPVAction} & \texttt{INT} & 1 & On bad PV: 0=Hold CV 1=Go Manual 2=Go Cfg\_IntlkCV \\
\texttt{Out\_CV} & \texttt{REAL} &  & Controller output (\%) \\
\texttt{Val\_SP} & \texttt{REAL} &  & Active setpoint after ramp and limits (EU) \\
\texttt{Val\_SPTarget} & \texttt{REAL} &  & Target setpoint before ramp (EU) \\
\texttt{Val\_PV} & \texttt{REAL} &  & Process value (EU) \\
\texttt{Val\_Err} & \texttt{REAL} &  & Error SP --- PV (EU) \\
\texttt{Sts\_Man} & \texttt{BOOL} &  & In Manual \\
\texttt{Sts\_Auto} & \texttt{BOOL} &  & In Auto \\
\texttt{Sts\_Cas} & \texttt{BOOL} &  & In Cascade \\
\texttt{Sts\_Track} & \texttt{BOOL} &  & Output tracking \\
\texttt{Sts\_Intlk} & \texttt{BOOL} &  & Interlock active \\
\texttt{Sts\_WindupHi} & \texttt{BOOL} &  & Output at high limit \\
\texttt{Sts\_WindupLo} & \texttt{BOOL} &  & Output at low limit \\
\texttt{Sts\_SPRamping} & \texttt{BOOL} &  & Setpoint ramping \\
\texttt{Sts\_Mode} & \texttt{INT} &  & 0=Man 1=Auto 2=Cas 3=Track 4=Intlk \\
\texttt{Alm\_DevH} & \texttt{T\_AlmLim} &  & High deviation (PV --- SP), EU \\
\texttt{Alm\_DevL} & \texttt{T\_AlmLim} &  & Low deviation (SP --- PV), EU \\
\texttt{Alm\_PVBad} & \texttt{T\_Alarm} &  & PV bad \\
\bottomrule
\end{longtable}
\end{small}

\section{Rotating Equipment}

\subsection{T\_PumpProt}\label{sec:pclib-t-pumpprot}

Pump protection layer: dry run, low flow, deadhead, seal, bearing.

\par\smallskip{\small \textbf{ISA-88 role:} Control module \quad \textbf{MTP:} - \quad \textbf{Standards:} ISA-88 /\allowbreak{} IEC 61512-1, ISA-18.2 /\allowbreak{} IEC 62682, ISA-101.01}\par

\begin{small}
\begin{longtable}{@{}P{3.70cm}P{2.20cm}P{1.70cm}P{8.91cm}@{}}
\caption{T\_PumpProt Members}\label{tab:pclib-t-pumpprot} \\
\toprule
\rowcolor{tblhead}\textbf{Member} & \textbf{Type} & \textbf{Init} & \textbf{Description} \\
\midrule
\endfirsthead
\multicolumn{4}{@{}l}{\small\textit{Table \thetable{} continued}} \\
\toprule
\rowcolor{tblhead}\textbf{Member} & \textbf{Type} & \textbf{Init} & \textbf{Description} \\
\midrule
\endhead
\texttt{Cfg\_Tag} & \texttt{STRING[20]} &  & Equipment /\allowbreak{} instrument tag, e.g. FT-101 \\
\texttt{Cfg\_Desc} & \texttt{STRING[40]} &  & Description shown on HMI faceplate \\
\texttt{Inp\_Running} & \texttt{BOOL} &  & Pump running (from Motor /\allowbreak{} VFD Sts\_Running) \\
\texttt{Inp\_SuctPress} & \texttt{REAL} &  & Suction pressure (EU) \\
\texttt{Inp\_DischPress} & \texttt{REAL} &  & Discharge pressure (EU) \\
\texttt{Inp\_Flow} & \texttt{REAL} &  & Pump flow (EU) \\
\texttt{Inp\_DryRun} & \texttt{BOOL} &  & Dry-run switch /\allowbreak{} low level active \\
\texttt{Inp\_SealLeak} & \texttt{BOOL} &  & Seal leak /\allowbreak{} seal pot alarm \\
\texttt{Inp\_BrgTemp} & \texttt{REAL} &  & Bearing temperature (degC) \\
\texttt{Inp\_Vib} & \texttt{REAL} &  & Vibration (mm/\allowbreak{}s) \\
\texttt{Cfg\_MinFlow} & \texttt{REAL} & 0.0 & Minimum continuous flow (EU) \\
\texttt{Cfg\_MinFlowDly} & \texttt{TIME} & T\#30S & Low-flow delay after start \\
\texttt{Cfg\_DeadheadDP} & \texttt{REAL} & 0.0 & Differential pressure indicating deadhead (EU) \\
\texttt{Cfg\_DryRunDly} & \texttt{TIME} & T\#5S & Dry-run delay \\
\texttt{Cfg\_TripDryRun} & \texttt{BOOL} & TRUE & Trip on dry run \\
\texttt{Cfg\_TripLowFlow} & \texttt{BOOL} &  & Trip on low flow \\
\texttt{Cfg\_TripDeadhead} & \texttt{BOOL} & TRUE & Trip on deadhead \\
\texttt{Cfg\_TripSeal} & \texttt{BOOL} &  & Trip on seal leak \\
\texttt{Cmd\_Reset} & \texttt{BOOL} &  & Reset latched trips (one-shot) \\
\texttt{Out\_IntlkOK} & \texttt{BOOL} &  & Protections healthy --- wire to Motor/\allowbreak{}VFD Inp\_IntlkOK \\
\texttt{Val\_DP} & \texttt{REAL} &  & Differential pressure (EU) \\
\texttt{Alm\_DryRun} & \texttt{T\_Alarm} &  & Dry run \\
\texttt{Alm\_LowFlow} & \texttt{T\_Alarm} &  & Low flow (below minimum continuous flow) \\
\texttt{Alm\_Deadhead} & \texttt{T\_Alarm} &  & Deadhead (high DP, no flow) \\
\texttt{Alm\_SealLeak} & \texttt{T\_Alarm} &  & Seal leak \\
\texttt{Alm\_BrgTemp} & \texttt{T\_AlmLim} &  & High bearing temperature \\
\texttt{Alm\_Vib} & \texttt{T\_AlmLim} &  & High vibration \\
\texttt{Alm\_LoSuct} & \texttt{T\_AlmLim} &  & Low suction pressure \\
\bottomrule
\end{longtable}
\end{small}

\subsection{T\_PumpDS}\label{sec:pclib-t-pumpds}

Duty/\allowbreak{}standby/\allowbreak{}lead-lag coordinator for up to 4 pumps.

\par\smallskip{\small \textbf{ISA-88 role:} Equipment module \quad \textbf{MTP:} - \quad \textbf{Standards:} ISA-88 /\allowbreak{} IEC 61512-1, ISA-18.2 /\allowbreak{} IEC 62682, ISA-101.01}\par

\begin{small}
\begin{longtable}{@{}P{3.70cm}P{2.20cm}P{1.70cm}P{8.91cm}@{}}
\caption{T\_PumpDS Members}\label{tab:pclib-t-pumpds} \\
\toprule
\rowcolor{tblhead}\textbf{Member} & \textbf{Type} & \textbf{Init} & \textbf{Description} \\
\midrule
\endfirsthead
\multicolumn{4}{@{}l}{\small\textit{Table \thetable{} continued}} \\
\toprule
\rowcolor{tblhead}\textbf{Member} & \textbf{Type} & \textbf{Init} & \textbf{Description} \\
\midrule
\endhead
\texttt{Cfg\_Tag} & \texttt{STRING[20]} &  & Equipment /\allowbreak{} instrument tag, e.g. FT-101 \\
\texttt{Cfg\_Desc} & \texttt{STRING[40]} &  & Description shown on HMI faceplate \\
\texttt{Mode} & \texttt{T\_Mode} &  & Mode block (Auto = Prog\_NumReq, Man = Cmd\_NumReq) \\
\texttt{Inp\_Avail} & \texttt{ARRAY[1..4] OF BOOL} &  & Pump n available (Motor Sts\_Avail) \\
\texttt{Inp\_Running} & \texttt{ARRAY[1..4] OF BOOL} &  & Pump n running \\
\texttt{Inp\_Fault} & \texttt{ARRAY[1..4] OF BOOL} &  & Pump n faulted \\
\texttt{Inp\_RunHrs} & \texttt{ARRAY[1..4] OF REAL} &  & Pump n run hours (for equalisation) \\
\texttt{Cfg\_NumPumps} & \texttt{INT} & 2 & Number of pumps in group (1-4) \\
\texttt{Cfg\_Strategy} & \texttt{INT} & 1 & Rotation: 0=Fixed lead 1=Rotate on stop 2=Equalise hours 3=Timed \\
\texttt{Cfg\_RotateHrs} & \texttt{REAL} & 168.0 & Timed rotation interval (h) \\
\texttt{Cfg\_StageDly} & \texttt{TIME} & T\#10S & Delay between staging pumps on/\allowbreak{}off \\
\texttt{Cfg\_ChgOnFault} & \texttt{BOOL} & TRUE & Automatic changeover to standby on fault \\
\texttt{Cmd\_NumReq} & \texttt{INT} & 0 & Operator: number of pumps required (Manual) \\
\texttt{Cmd\_Lead} & \texttt{INT} & 1 & Operator-selected lead pump (Fixed strategy) \\
\texttt{Cmd\_Rotate} & \texttt{BOOL} &  & Force rotation now (one-shot) \\
\texttt{Prog\_NumReq} & \texttt{INT} & 0 & Program: number of pumps required (Auto) \\
\texttt{Out\_RunReq} & \texttt{ARRAY[1..4] OF BOOL} &  & Run request to pump n (wire to Prog\_Start) \\
\texttt{Sts\_Order} & \texttt{ARRAY[1..4] OF INT} &  & Rank of pump n: 1=Lead 2=Lag1 ... \\
\texttt{Sts\_Lead} & \texttt{INT} &  & Current lead pump \\
\texttt{Sts\_NumRunning} & \texttt{INT} &  & Pumps running \\
\texttt{Sts\_NumAvail} & \texttt{INT} &  & Pumps available \\
\texttt{Val\_HrsToRotate} & \texttt{REAL} &  & Hours until next timed rotation \\
\texttt{Alm\_NoStandby} & \texttt{T\_Alarm} &  & No standby pump available \\
\texttt{Alm\_NoneAvail} & \texttt{T\_Alarm} &  & Demand present but no pump available \\
\texttt{Sts\_ChgOver} & \texttt{BOOL} &  & Automatic changeover occurred (event, not an alarm --- ISA-18.2) \\
\texttt{Val\_ChgOverCount} & \texttt{UDINT} &  & Automatic changeover counter \\
\bottomrule
\end{longtable}
\end{small}

\subsection{T\_Agitator}\label{sec:pclib-t-agitator}

Agitator /\allowbreak{} mixer: VFD + submergence, seal and gearbox protection.

\par\smallskip{\small \textbf{ISA-88 role:} Equipment module \quad \textbf{MTP:} MonAnaDrv (drive) \quad \textbf{Standards:} ISA-88 /\allowbreak{} IEC 61512-1, ISA-18.2 /\allowbreak{} IEC 62682, ISA-101.01}\par

\begin{small}
\begin{longtable}{@{}P{3.70cm}P{2.20cm}P{1.70cm}P{8.91cm}@{}}
\caption{T\_Agitator Members}\label{tab:pclib-t-agitator} \\
\toprule
\rowcolor{tblhead}\textbf{Member} & \textbf{Type} & \textbf{Init} & \textbf{Description} \\
\midrule
\endfirsthead
\multicolumn{4}{@{}l}{\small\textit{Table \thetable{} continued}} \\
\toprule
\rowcolor{tblhead}\textbf{Member} & \textbf{Type} & \textbf{Init} & \textbf{Description} \\
\midrule
\endhead
\texttt{Drv} & \texttt{T\_VFD} &  & Drive (fixed-speed agitators: ignore speed members) \\
\texttt{Inp\_Level} & \texttt{REAL} &  & Vessel level (EU) \\
\texttt{Inp\_SealOK} & \texttt{BOOL} &  & Seal /\allowbreak{} seal support system healthy \\
\texttt{Inp\_GbxTemp} & \texttt{REAL} &  & Gearbox oil temperature (degC) \\
\texttt{Inp\_Vib} & \texttt{REAL} &  & Vibration (mm/\allowbreak{}s) \\
\texttt{Cfg\_MinLevel} & \texttt{REAL} & 0.0 & Minimum level to run (impeller submergence) \\
\texttt{Cfg\_LvlHyst} & \texttt{REAL} & 0.0 & Level hysteresis (EU) \\
\texttt{Cfg\_TripLowLvl} & \texttt{BOOL} & TRUE & Trip running agitator on low level \\
\texttt{Out\_IntlkOK} & \texttt{BOOL} &  & Wire to Drv.Inp\_IntlkOK \\
\texttt{Out\_PermOK} & \texttt{BOOL} &  & Wire to Drv.Inp\_PermOK \\
\texttt{Sts\_LvlOK} & \texttt{BOOL} &  & Level above minimum \\
\texttt{Alm\_LowLvlRun} & \texttt{T\_Alarm} &  & Running /\allowbreak{} start requested below minimum level \\
\texttt{Alm\_Seal} & \texttt{T\_Alarm} &  & Seal fault \\
\texttt{Alm\_GbxTemp} & \texttt{T\_AlmLim} &  & High gearbox temperature \\
\texttt{Alm\_Vib} & \texttt{T\_AlmLim} &  & High vibration \\
\bottomrule
\end{longtable}
\end{small}

\subsection{T\_Compressor}\label{sec:pclib-t-compressor}

Compressor /\allowbreak{} blower: load-unload, anti-cycle, lube and thermal protection.

\par\smallskip{\small \textbf{ISA-88 role:} Equipment module \quad \textbf{MTP:} MonAnaDrv (drive) \quad \textbf{Standards:} ISA-88 /\allowbreak{} IEC 61512-1, ISA-18.2 /\allowbreak{} IEC 62682, ISA-101.01}\par

\begin{small}
\begin{longtable}{@{}P{3.70cm}P{2.20cm}P{1.70cm}P{8.91cm}@{}}
\caption{T\_Compressor Members}\label{tab:pclib-t-compressor} \\
\toprule
\rowcolor{tblhead}\textbf{Member} & \textbf{Type} & \textbf{Init} & \textbf{Description} \\
\midrule
\endfirsthead
\multicolumn{4}{@{}l}{\small\textit{Table \thetable{} continued}} \\
\toprule
\rowcolor{tblhead}\textbf{Member} & \textbf{Type} & \textbf{Init} & \textbf{Description} \\
\midrule
\endhead
\texttt{Drv} & \texttt{T\_VFD} &  & Drive (fixed-speed: ignore speed members) \\
\texttt{Inp\_SuctPress} & \texttt{REAL} &  & Suction pressure (EU) \\
\texttt{Inp\_DischPress} & \texttt{REAL} &  & Discharge /\allowbreak{} header pressure (EU) \\
\texttt{Inp\_DischTemp} & \texttt{REAL} &  & Discharge temperature (degC) \\
\texttt{Inp\_OilPress} & \texttt{REAL} &  & Lube oil pressure (EU) \\
\texttt{Inp\_OilTemp} & \texttt{REAL} &  & Lube oil temperature (degC) \\
\texttt{Inp\_Vib} & \texttt{REAL} &  & Vibration (mm/\allowbreak{}s) \\
\texttt{Inp\_LoadFbk} & \texttt{BOOL} &  & Loaded feedback (inlet valve open) \\
\texttt{Cfg\_LoadPress} & \texttt{REAL} & 0.0 & Load (cut-in) pressure \\
\texttt{Cfg\_UnloadPress} & \texttt{REAL} & 0.0 & Unload (cut-out) pressure \\
\texttt{Cfg\_MinRun} & \texttt{TIME} & T\#5M & Minimum run time \\
\texttt{Cfg\_MinStop} & \texttt{TIME} & T\#2M & Minimum stop time (anti-cycle) \\
\texttt{Cfg\_IdleTime} & \texttt{TIME} & T\#10M & Unloaded run time before auto stop \\
\texttt{Cfg\_OilBypass} & \texttt{TIME} & T\#10S & Low oil pressure bypass after start \\
\texttt{Out\_Load} & \texttt{BOOL} &  & Load solenoid /\allowbreak{} inlet valve \\
\texttt{Out\_BlowOff} & \texttt{BOOL} &  & Blow-off /\allowbreak{} recycle valve open \\
\texttt{Out\_IntlkOK} & \texttt{BOOL} &  & Wire to Drv.Inp\_IntlkOK \\
\texttt{Sts\_Loaded} & \texttt{BOOL} &  & Loaded \\
\texttt{Sts\_Idle} & \texttt{BOOL} &  & Running unloaded \\
\texttt{Sts\_AntiCycle} & \texttt{BOOL} &  & Waiting on minimum stop time \\
\texttt{Sts\_State} & \texttt{INT} &  & 0=Off 1=Starting 2=Unloaded 3=Loaded 4=Idle stop 5=Anti-cycle 6=Tripped \\
\texttt{Val\_TimeToStart} & \texttt{TIME} &  & Anti-cycle time remaining \\
\texttt{Val\_PressRatio} & \texttt{REAL} &  & Compression ratio (absolute) \\
\texttt{Alm\_HiDischPress} & \texttt{T\_AlmLim} &  & High discharge pressure \\
\texttt{Alm\_HiDischTemp} & \texttt{T\_AlmLim} &  & High discharge temperature \\
\texttt{Alm\_LoOilPress} & \texttt{T\_AlmLim} &  & Low lube oil pressure \\
\texttt{Alm\_HiOilTemp} & \texttt{T\_AlmLim} &  & High lube oil temperature \\
\texttt{Alm\_Vib} & \texttt{T\_AlmLim} &  & High vibration \\
\texttt{Alm\_LoSuct} & \texttt{T\_AlmLim} &  & Low suction pressure \\
\bottomrule
\end{longtable}
\end{small}

\subsection{T\_Conveyor}\label{sec:pclib-t-conveyor}

Belt /\allowbreak{} screw conveyor: start warning, zero-speed, pull-cord, alignment.

\par\smallskip{\small \textbf{ISA-88 role:} Equipment module \quad \textbf{MTP:} MonBinDrv (drive) \quad \textbf{Standards:} ISA-88 /\allowbreak{} IEC 61512-1, ISA-18.2 /\allowbreak{} IEC 62682, ISA-101.01}\par

\begin{small}
\begin{longtable}{@{}P{3.70cm}P{2.20cm}P{1.70cm}P{8.91cm}@{}}
\caption{T\_Conveyor Members}\label{tab:pclib-t-conveyor} \\
\toprule
\rowcolor{tblhead}\textbf{Member} & \textbf{Type} & \textbf{Init} & \textbf{Description} \\
\midrule
\endfirsthead
\multicolumn{4}{@{}l}{\small\textit{Table \thetable{} continued}} \\
\toprule
\rowcolor{tblhead}\textbf{Member} & \textbf{Type} & \textbf{Init} & \textbf{Description} \\
\midrule
\endhead
\texttt{Mtr} & \texttt{T\_Motor} &  & Drive motor \\
\texttt{Inp\_MotionOK} & \texttt{BOOL} &  & Zero-speed switch: motion detected \\
\texttt{Inp\_PullCordOK} & \texttt{BOOL} &  & Pull-cord /\allowbreak{} lanyard switches healthy \\
\texttt{Inp\_AlignLOK} & \texttt{BOOL} &  & Belt alignment left healthy \\
\texttt{Inp\_AlignROK} & \texttt{BOOL} &  & Belt alignment right healthy \\
\texttt{Inp\_ChuteOK} & \texttt{BOOL} &  & Chute not blocked \\
\texttt{Inp\_BeltRipOK} & \texttt{BOOL} &  & Belt rip detector healthy \\
\texttt{Inp\_DownstrRun} & \texttt{BOOL} &  & Downstream conveyor running \\
\texttt{Inp\_Rate} & \texttt{REAL} &  & Belt scale rate (t/\allowbreak{}h), optional \\
\texttt{Cfg\_UseDownstr} & \texttt{BOOL} &  & Require downstream running (cascade start/\allowbreak{}trip) \\
\texttt{Cfg\_WarnTime} & \texttt{TIME} & T\#10S & Pre-start warning horn time \\
\texttt{Cfg\_MotionDly} & \texttt{TIME} & T\#5S & Time after start to detect motion \\
\texttt{Cfg\_PurgeTime} & \texttt{TIME} & T\#0S & Run-out time before stopping (empty belt) \\
\texttt{Out\_Horn} & \texttt{BOOL} &  & Start warning horn \\
\texttt{Out\_Beacon} & \texttt{BOOL} &  & Start warning beacon \\
\texttt{Out\_IntlkOK} & \texttt{BOOL} &  & Wire to Mtr.Inp\_IntlkOK \\
\texttt{Out\_PermOK} & \texttt{BOOL} &  & Wire to Mtr.Inp\_PermOK \\
\texttt{Out\_StartReq} & \texttt{BOOL} &  & Start request to Mtr after warning completes \\
\texttt{Sts\_Warning} & \texttt{BOOL} &  & Pre-start warning in progress \\
\texttt{Sts\_Purging} & \texttt{BOOL} &  & Running out before stop \\
\texttt{Sts\_State} & \texttt{INT} &  & 0=Stopped 1=Warning 2=Starting 3=Running 4=Purging 5=Tripped \\
\texttt{Val\_Total} & \texttt{LREAL} &  & Total conveyed (t) \\
\texttt{Alm\_ZeroSpd} & \texttt{T\_Alarm} &  & Zero-speed /\allowbreak{} belt slip \\
\texttt{Alm\_PullCord} & \texttt{T\_Alarm} &  & Pull-cord operated \\
\texttt{Alm\_Misalign} & \texttt{T\_Alarm} &  & Belt misalignment \\
\texttt{Alm\_Chute} & \texttt{T\_Alarm} &  & Chute blocked \\
\texttt{Alm\_BeltRip} & \texttt{T\_Alarm} &  & Belt rip \\
\texttt{Alm\_DownstrStop} & \texttt{T\_Alarm} &  & Tripped by downstream conveyor \\
\bottomrule
\end{longtable}
\end{small}

\subsection{T\_DiaComp}\label{sec:pclib-t-diacomp}

Diaphragm compressor unit: one motor, two diaphragm stages, leak detection, hydraulic oil, purge.

\par\smallskip{\small \textbf{ISA-88 role:} Equipment module \quad \textbf{MTP:} MonBinDrv (drive) \quad \textbf{Standards:} ISA-88 /\allowbreak{} IEC 61512-1, ISA-5.1, ISA-18.2 /\allowbreak{} IEC 62682, ISA-101.01}\par

Built for the PDC/Fluitron-style two-stage diaphragm compressor (Section~\ref{sec:pdc-compressor}): one motor turns the crankshaft for both stages, so the unit owns a single nested \texttt{T\_Motor} and drives it through \texttt{Out\_StartReq}, \texttt{Out\_StopReq}, \texttt{Out\_IntlkOK}, and \texttt{Out\_PermOK}. \texttt{Sts\_FirstOut} records the first trip cause for the HMI. Trip setpoints and shutdown actions come from the vendor documentation and the hazard review, not from the initial values below.

\begin{small}
\begin{longtable}{@{}P{3.70cm}P{2.20cm}P{1.70cm}P{8.91cm}@{}}
\caption{T\_DiaComp Members}\label{tab:pclib-t-diacomp} \\
\toprule
\rowcolor{tblhead}\textbf{Member} & \textbf{Type} & \textbf{Init} & \textbf{Description} \\
\midrule
\endfirsthead
\multicolumn{4}{@{}l}{\small\textit{Table \thetable{} continued}} \\
\toprule
\rowcolor{tblhead}\textbf{Member} & \textbf{Type} & \textbf{Init} & \textbf{Description} \\
\midrule
\endhead
\texttt{Cfg\_Tag} & \texttt{STRING[20]} &  & Equipment /\allowbreak{} instrument tag, e.g. FT-101 \\
\texttt{Cfg\_Desc} & \texttt{STRING[40]} &  & Description shown on HMI faceplate \\
\texttt{Mode} & \texttt{T\_Mode} &  & Unit mode block (Auto = Prog\_ commands, Man = operator) \\
\texttt{Mtr} & \texttt{T\_Motor} &  & Drive motor (one motor drives both stages through the crankshaft) \\
\texttt{Cfg\_PressEU} & \texttt{STRING[8]} &  & Pressure units text, e.g. barg, psig \\
\texttt{Inp\_SuctPress} & \texttt{REAL} &  & Stage 1 suction pressure (Cfg\_PressEU) \\
\texttt{Inp\_InterPress} & \texttt{REAL} &  & Interstage pressure: stage 1 discharge /\allowbreak{} stage 2 suction \\
\texttt{Inp\_DischPress} & \texttt{REAL} &  & Stage 2 (final) discharge pressure \\
\texttt{Inp\_S1DischTemp} & \texttt{REAL} &  & Stage 1 discharge gas temperature (degC) \\
\texttt{Inp\_S2DischTemp} & \texttt{REAL} &  & Stage 2 discharge gas temperature (degC) \\
\texttt{Inp\_S1OilPress} & \texttt{REAL} &  & Stage 1 head hydraulic oil pressure (Cfg\_PressEU) \\
\texttt{Inp\_S2OilPress} & \texttt{REAL} &  & Stage 2 head hydraulic oil pressure (Cfg\_PressEU) \\
\texttt{Inp\_CaseOilPress} & \texttt{REAL} &  & Crankcase lube oil pressure (Cfg\_PressEU) \\
\texttt{Inp\_CaseOilTemp} & \texttt{REAL} &  & Crankcase lube oil temperature (degC) \\
\texttt{Inp\_CWTemp} & \texttt{REAL} &  & Cooling water return temperature (degC) \\
\texttt{Inp\_S1LeakOK} & \texttt{BOOL} &  & Stage 1 diaphragm leak detection healthy (no leak between diaphragms) \\
\texttt{Inp\_S2LeakOK} & \texttt{BOOL} &  & Stage 2 diaphragm leak detection healthy (no leak between diaphragms) \\
\texttt{Inp\_CWFlowOK} & \texttt{BOOL} &  & Cooling water flow switch healthy \\
\texttt{Inp\_GasDetOK} & \texttt{BOOL} &  & Area gas detection healthy (from T\_GasDet /\allowbreak{} ESD) \\
\texttt{Inp\_ExtIntlkOK} & \texttt{BOOL} &  & External interlocks healthy (ESD, upstream /\allowbreak{} downstream) \\
\texttt{Inp\_SuctVlvOpen} & \texttt{BOOL} &  & Suction isolation valve confirmed open \\
\texttt{Inp\_DischVlvOpen} & \texttt{BOOL} &  & Discharge isolation valve confirmed open \\
\texttt{Cfg\_CtrlMode} & \texttt{INT} & 0 & 0=Start/\allowbreak{}stop by command only 1=Auto start/\allowbreak{}stop on Inp\_DischPress 2=Run continuous, recycle on Inp\_DischPress \\
\texttt{Cfg\_LoadPress} & \texttt{REAL} & 0.0 & Load /\allowbreak{} start (cut-in) discharge pressure \\
\texttt{Cfg\_UnloadPress} & \texttt{REAL} & 0.0 & Unload /\allowbreak{} stop (cut-out) discharge pressure \\
\texttt{Cfg\_UnloadedStart} & \texttt{TIME} & T\#10S & Run on recycle after start before loading \\
\texttt{Cfg\_OilBypass} & \texttt{TIME} & T\#10S & Low oil pressure trips bypassed after start \\
\texttt{Cfg\_MinRun} & \texttt{TIME} & T\#5M & Minimum run time \\
\texttt{Cfg\_MinStop} & \texttt{TIME} & T\#2M & Minimum stop time (anti-cycle) \\
\texttt{Cfg\_IdleTime} & \texttt{TIME} & T\#10M & Unloaded run time before auto stop (Cfg\_CtrlMode 1) \\
\texttt{Cfg\_LeakDly} & \texttt{TIME} & T\#2S & Diaphragm leak confirmation delay \\
\texttt{Cfg\_TripOnLeak} & \texttt{BOOL} & TRUE & Trip on diaphragm leak (FALSE = alarm only) \\
\texttt{Cfg\_PurgeOnLeak} & \texttt{BOOL} &  & Start the purge automatically after a leak trip \\
\texttt{Cfg\_PurgeTime} & \texttt{TIME} & T\#5M & Purge duration \\
\texttt{Cfg\_DiaHrs} & \texttt{REAL} & 0.0 & Diaphragm replacement interval (run hours), 0 = off \\
\texttt{Cfg\_PAtm} & \texttt{REAL} & 0.0 & Added to gauge pressures for compression ratio; 0 if transmitters read absolute \\
\texttt{Cmd\_Start} & \texttt{BOOL} &  & Operator start (one-shot) \\
\texttt{Cmd\_Stop} & \texttt{BOOL} &  & Operator stop (one-shot) \\
\texttt{Cmd\_Reset} & \texttt{BOOL} &  & Reset unit trips (one-shot) \\
\texttt{Cmd\_Load} & \texttt{BOOL} &  & Operator load: close recycle (one-shot, Manual) \\
\texttt{Cmd\_Unload} & \texttt{BOOL} &  & Operator unload: open recycle (one-shot, Manual) \\
\texttt{Cmd\_Purge} & \texttt{BOOL} &  & Start purge (one-shot, unit stopped) \\
\texttt{Cmd\_PurgeStop} & \texttt{BOOL} &  & Abort purge (one-shot) \\
\texttt{Cmd\_ResetDiaS1} & \texttt{BOOL} &  & Stage 1 diaphragms replaced: reset hours (one-shot) \\
\texttt{Cmd\_ResetDiaS2} & \texttt{BOOL} &  & Stage 2 diaphragms replaced: reset hours (one-shot) \\
\texttt{Prog\_Start} & \texttt{BOOL} &  & Program /\allowbreak{} sequence start request (Auto) \\
\texttt{Prog\_Stop} & \texttt{BOOL} &  & Program /\allowbreak{} sequence stop request (Auto) \\
\texttt{Out\_StartReq} & \texttt{BOOL} &  & Start request --- wire to Mtr.Prog\_Start \\
\texttt{Out\_StopReq} & \texttt{BOOL} &  & Stop request --- wire to Mtr.Prog\_Stop \\
\texttt{Out\_IntlkOK} & \texttt{BOOL} &  & Unit trips healthy --- wire to Mtr.Inp\_IntlkOK \\
\texttt{Out\_PermOK} & \texttt{BOOL} &  & Start permissives healthy --- wire to Mtr.Inp\_PermOK \\
\texttt{Out\_Recycle} & \texttt{BOOL} &  & Recycle /\allowbreak{} bypass valve open (unloaded start, unload) \\
\texttt{Out\_SuctVlv} & \texttt{BOOL} &  & Open suction isolation valve \\
\texttt{Out\_DischVlv} & \texttt{BOOL} &  & Open discharge isolation valve \\
\texttt{Out\_CWVlv} & \texttt{BOOL} &  & Open cooling water valve \\
\texttt{Out\_PurgeVlv} & \texttt{BOOL} &  & Open purge gas valve \\
\texttt{Sim\_En} & \texttt{BOOL} &  & Simulation: ignore field I/\allowbreak{}O, logic generates feedback \\
\texttt{Sts\_State} & \texttt{INT} &  & 0=Stopped 1=Pre-start 2=Starting unloaded 3=Loaded 4=Unloaded (recycle) 5=Stopping 6=Anti-cycle 7=Purging 8=Tripped 9=Not available \\
\texttt{Sts\_Loaded} & \texttt{BOOL} &  & Running loaded \\
\texttt{Sts\_Unloaded} & \texttt{BOOL} &  & Running on recycle \\
\texttt{Sts\_AntiCycle} & \texttt{BOOL} &  & Waiting on minimum stop time \\
\texttt{Sts\_Purging} & \texttt{BOOL} &  & Purge in progress \\
\texttt{Sts\_Tripped} & \texttt{BOOL} &  & Unit tripped --- reset required \\
\texttt{Sts\_ResetReq} & \texttt{BOOL} &  & Trip causes cleared, reset required \\
\texttt{Sts\_FirstOut} & \texttt{INT} &  & First trip cause: 0=None 1=S1 leak 2=S2 leak 3=Case oil 4=S1 head oil 5=S2 head oil 6=S1 disch temp 7=S2 disch temp 8=Interstage press 9=Disch press 10=Suction press 11=Cooling water 12=Gas detection 13=External interlock 14=Motor \\
\texttt{Sts\_DiaDueS1} & \texttt{BOOL} &  & Stage 1 diaphragm replacement due \\
\texttt{Sts\_DiaDueS2} & \texttt{BOOL} &  & Stage 2 diaphragm replacement due \\
\texttt{Val\_TimeToStart} & \texttt{TIME} &  & Anti-cycle time remaining \\
\texttt{Val\_PurgeLeft} & \texttt{TIME} &  & Purge time remaining \\
\texttt{Val\_S1Ratio} & \texttt{REAL} &  & Stage 1 compression ratio (absolute) \\
\texttt{Val\_S2Ratio} & \texttt{REAL} &  & Stage 2 compression ratio (absolute) \\
\texttt{Val\_S1DiaHrs} & \texttt{REAL} &  & Stage 1 run hours since diaphragm change \\
\texttt{Val\_S2DiaHrs} & \texttt{REAL} &  & Stage 2 run hours since diaphragm change \\
\texttt{Alm\_S1Leak} & \texttt{T\_Alarm} &  & Stage 1 diaphragm leak \\
\texttt{Alm\_S2Leak} & \texttt{T\_Alarm} &  & Stage 2 diaphragm leak \\
\texttt{Alm\_LoCaseOil} & \texttt{T\_AlmLim} &  & Low crankcase oil pressure \\
\texttt{Alm\_HiCaseOilT} & \texttt{T\_AlmLim} &  & High crankcase oil temperature \\
\texttt{Alm\_LoS1Oil} & \texttt{T\_AlmLim} &  & Low stage 1 head hydraulic oil pressure \\
\texttt{Alm\_LoS2Oil} & \texttt{T\_AlmLim} &  & Low stage 2 head hydraulic oil pressure \\
\texttt{Alm\_HiS1DischT} & \texttt{T\_AlmLim} &  & High stage 1 discharge temperature \\
\texttt{Alm\_HiS2DischT} & \texttt{T\_AlmLim} &  & High stage 2 discharge temperature \\
\texttt{Alm\_HiInterPress} & \texttt{T\_AlmLim} &  & High interstage pressure \\
\texttt{Alm\_HiDischPress} & \texttt{T\_AlmLim} &  & High final discharge pressure \\
\texttt{Alm\_LoSuct} & \texttt{T\_AlmLim} &  & Low suction pressure \\
\texttt{Alm\_HiCWTemp} & \texttt{T\_AlmLim} &  & High cooling water return temperature \\
\texttt{Alm\_CWFail} & \texttt{T\_Alarm} &  & Cooling water flow failure \\
\texttt{Alm\_GasDet} & \texttt{T\_Alarm} &  & Unit tripped by area gas detection \\
\texttt{Alm\_ExtIntlk} & \texttt{T\_Alarm} &  & Unit tripped by external interlock \\
\texttt{Alm\_DiaDue} & \texttt{T\_Alarm} &  & Diaphragm replacement due (either stage) \\
\bottomrule
\end{longtable}
\end{small}

\section{Process Equipment}

\subsection{T\_Tank}\label{sec:pclib-t-tank}

Tank /\allowbreak{} vessel: level-to-volume, inventory, fill/\allowbreak{}drain permissives, alarms.

\par\smallskip{\small \textbf{ISA-88 role:} Equipment module /\allowbreak{} Unit vessel \quad \textbf{MTP:} - \quad \textbf{Standards:} ISA-88 /\allowbreak{} IEC 61512-1, ISA-5.1, ISA-18.2 /\allowbreak{} IEC 62682, ISA-101.01}\par

\begin{small}
\begin{longtable}{@{}P{3.70cm}P{2.20cm}P{1.70cm}P{8.91cm}@{}}
\caption{T\_Tank Members}\label{tab:pclib-t-tank} \\
\toprule
\rowcolor{tblhead}\textbf{Member} & \textbf{Type} & \textbf{Init} & \textbf{Description} \\
\midrule
\endfirsthead
\multicolumn{4}{@{}l}{\small\textit{Table \thetable{} continued}} \\
\toprule
\rowcolor{tblhead}\textbf{Member} & \textbf{Type} & \textbf{Init} & \textbf{Description} \\
\midrule
\endhead
\texttt{Cfg\_Tag} & \texttt{STRING[20]} &  & Equipment /\allowbreak{} instrument tag, e.g. FT-101 \\
\texttt{Cfg\_Desc} & \texttt{STRING[40]} &  & Description shown on HMI faceplate \\
\texttt{Cfg\_Product} & \texttt{STRING[24]} &  & Stored product \\
\texttt{Cfg\_LvlEU} & \texttt{STRING[8]} &  & Level units text \\
\texttt{Cfg\_VolEU} & \texttt{STRING[8]} &  & Volume units text \\
\texttt{Inp\_Level} & \texttt{REAL} &  & Level (Cfg\_LvlEU) from AI Val\_PV \\
\texttt{Inp\_LevelBad} & \texttt{BOOL} &  & Level signal bad \\
\texttt{Inp\_LSHH} & \texttt{BOOL} &  & High-high level switch active \\
\texttt{Inp\_LSLL} & \texttt{BOOL} &  & Low-low level switch active \\
\texttt{Inp\_Temp} & \texttt{REAL} &  & Temperature (degC) \\
\texttt{Inp\_Press} & \texttt{REAL} &  & Pressure (EU) \\
\texttt{Cfg\_MaxLevel} & \texttt{REAL} & 100.0 & Level at 100\% full \\
\texttt{Cfg\_Capacity} & \texttt{REAL} & 100.0 & Volume at 100\% full \\
\texttt{Cfg\_Shape} & \texttt{INT} & 0 & 0=Vertical /\allowbreak{} linear 1=Horizontal cylinder 2=Strapping table \\
\texttt{Cfg\_Strap} & \texttt{ARRAY[0..10] OF REAL} &  & Strapping table: volume at 0,10..100\% level \\
\texttt{Cfg\_SG} & \texttt{REAL} & 1.0 & Specific gravity (for mass) \\
\texttt{Cfg\_FillStop} & \texttt{REAL} & 90.0 & Fill permissive removed above this level \\
\texttt{Cfg\_DrainStop} & \texttt{REAL} & 5.0 & Drain permissive removed below this level \\
\texttt{Cfg\_RateWindow} & \texttt{TIME} & T\#1M & Window for net-flow (rate of change) calc \\
\texttt{Prog\_EqState} & \texttt{INT} &  & ISA-88 equipment state request from unit /\allowbreak{} phase logic \\
\texttt{Out\_FillPerm} & \texttt{BOOL} &  & Filling permitted (wire to inlet pump /\allowbreak{} valve) \\
\texttt{Out\_DrainPerm} & \texttt{BOOL} &  & Draining permitted (wire to outlet pump /\allowbreak{} valve) \\
\texttt{Val\_Volume} & \texttt{REAL} &  & Volume (Cfg\_VolEU) \\
\texttt{Val\_PctFull} & \texttt{REAL} &  & Volume (\% of capacity) \\
\texttt{Val\_Ullage} & \texttt{REAL} &  & Remaining capacity (Cfg\_VolEU) \\
\texttt{Val\_Mass} & \texttt{REAL} &  & Mass (t, from volume x SG) \\
\texttt{Val\_NetFlow} & \texttt{REAL} &  & Net in/\allowbreak{}out flow (Cfg\_VolEU per min, + = filling) \\
\texttt{Val\_MinToFull} & \texttt{REAL} &  & Minutes to full at current rate \\
\texttt{Val\_MinToEmpty} & \texttt{REAL} &  & Minutes to empty at current rate \\
\texttt{Sts\_Filling} & \texttt{BOOL} &  & Level rising \\
\texttt{Sts\_Draining} & \texttt{BOOL} &  & Level falling \\
\texttt{Sts\_EqState} & \texttt{INT} &  & ISA-88 equipment state: 0=Available-clean 1=In use 2=Dirty 3=Cleaning (CIP) 4=Sterilising (SIP) 5=Maintenance 6=Out of service \\
\texttt{Alm\_HH} & \texttt{T\_AlmLim} &  & High-high level \\
\texttt{Alm\_H} & \texttt{T\_AlmLim} &  & High level \\
\texttt{Alm\_L} & \texttt{T\_AlmLim} &  & Low level \\
\texttt{Alm\_LL} & \texttt{T\_AlmLim} &  & Low-low level \\
\texttt{Alm\_LSHH} & \texttt{T\_Alarm} &  & High-high level switch \\
\texttt{Alm\_LSLL} & \texttt{T\_Alarm} &  & Low-low level switch \\
\texttt{Alm\_LvlDisc} & \texttt{T\_Alarm} &  & Analog level disagrees with level switch \\
\texttt{Alm\_HiTemp} & \texttt{T\_AlmLim} &  & High temperature \\
\texttt{Alm\_HiPress} & \texttt{T\_AlmLim} &  & High pressure \\
\bottomrule
\end{longtable}
\end{small}

\subsection{T\_Totalizer}\label{sec:pclib-t-totalizer}

Flow totalizer + batch delivery: analog or pulse, day/\allowbreak{}hour totals.

\par\smallskip{\small \textbf{ISA-88 role:} Control module \quad \textbf{MTP:} - \quad \textbf{Standards:} ISA-88 /\allowbreak{} IEC 61512-1, ISA-5.1, ISA-18.2 /\allowbreak{} IEC 62682, ISA-101.01}\par

\begin{small}
\begin{longtable}{@{}P{3.70cm}P{2.20cm}P{1.70cm}P{8.91cm}@{}}
\caption{T\_Totalizer Members}\label{tab:pclib-t-totalizer} \\
\toprule
\rowcolor{tblhead}\textbf{Member} & \textbf{Type} & \textbf{Init} & \textbf{Description} \\
\midrule
\endfirsthead
\multicolumn{4}{@{}l}{\small\textit{Table \thetable{} continued}} \\
\toprule
\rowcolor{tblhead}\textbf{Member} & \textbf{Type} & \textbf{Init} & \textbf{Description} \\
\midrule
\endhead
\texttt{Cfg\_Tag} & \texttt{STRING[20]} &  & Equipment /\allowbreak{} instrument tag, e.g. FT-101 \\
\texttt{Cfg\_Desc} & \texttt{STRING[40]} &  & Description shown on HMI faceplate \\
\texttt{Cfg\_RateEU} & \texttt{STRING[12]} &  & Rate units text, e.g. m3/\allowbreak{}h \\
\texttt{Cfg\_TotEU} & \texttt{STRING[12]} &  & Total units text, e.g. m3 \\
\texttt{Inp\_Rate} & \texttt{REAL} &  & Flow rate (Cfg\_RateEU) \\
\texttt{Inp\_RateBad} & \texttt{BOOL} &  & Rate signal bad \\
\texttt{Inp\_Pulses} & \texttt{UDINT} &  & Raw pulse counter (pulse mode) \\
\texttt{Cfg\_UsePulses} & \texttt{BOOL} &  & Totalise from pulses instead of rate \\
\texttt{Cfg\_PulseWt} & \texttt{REAL} & 1.0 & Total units per pulse \\
\texttt{Cfg\_TimeBase} & \texttt{REAL} & 3600.0 & Seconds per rate time unit: 1=/\allowbreak{}s 60=/\allowbreak{}min 3600=/\allowbreak{}h \\
\texttt{Cfg\_Conv} & \texttt{REAL} & 1.0 & Conversion rate units -\textgreater{} total units \\
\texttt{Cfg\_LoCutoff} & \texttt{REAL} & 0.0 & Rates below this are not totalised \\
\texttt{Cfg\_BadAction} & \texttt{INT} & 0 & On bad rate: 0=Stop totalising 1=Use last good rate \\
\texttt{Cfg\_DayStartHr} & \texttt{INT} & 0 & Hour the production day starts (0-23) \\
\texttt{Cfg\_BatchSP} & \texttt{REAL} & 0.0 & Batch target quantity \\
\texttt{Cfg\_PreAct} & \texttt{REAL} & 0.0 & Pre-act (in-flight) quantity \\
\texttt{Cfg\_AutoPreAct} & \texttt{BOOL} &  & Learn pre-act from overrun \\
\texttt{Cfg\_SlowPct} & \texttt{REAL} & 90.0 & Dribble /\allowbreak{} slow-fill at this \% of batch \\
\texttt{Cmd\_Start} & \texttt{BOOL} &  & Operator batch start (one-shot) \\
\texttt{Cmd\_Hold} & \texttt{BOOL} &  & Operator batch hold (one-shot) \\
\texttt{Cmd\_Reset} & \texttt{BOOL} &  & Operator batch reset (one-shot) \\
\texttt{Prog\_Start} & \texttt{BOOL} &  & Program batch start \\
\texttt{Prog\_Hold} & \texttt{BOOL} &  & Program batch hold \\
\texttt{Prog\_Reset} & \texttt{BOOL} &  & Program batch reset \\
\texttt{Out\_Slow} & \texttt{BOOL} &  & Slow-fill /\allowbreak{} dribble \\
\texttt{Out\_BatchDone} & \texttt{BOOL} &  & Batch target reached (close valve) \\
\texttt{Val\_Batch} & \texttt{LREAL} &  & Batch total \\
\texttt{Val\_Remaining} & \texttt{REAL} &  & Batch remaining \\
\texttt{Val\_Overrun} & \texttt{REAL} &  & Last batch overrun \\
\texttt{Val\_GrandTotal} & \texttt{LREAL} &  & Grand total (non-resettable) \\
\texttt{Val\_ThisHr} & \texttt{LREAL} &  & Total this hour \\
\texttt{Val\_LastHr} & \texttt{LREAL} &  & Total last hour \\
\texttt{Val\_Today} & \texttt{LREAL} &  & Total this production day \\
\texttt{Val\_Yday} & \texttt{LREAL} &  & Total previous production day \\
\texttt{Sts\_Running} & \texttt{BOOL} &  & Batch running \\
\texttt{Sts\_Hold} & \texttt{BOOL} &  & Batch held \\
\texttt{Sts\_Done} & \texttt{BOOL} &  & Batch complete \\
\texttt{Alm\_HiRate} & \texttt{T\_AlmLim} &  & High flow rate \\
\texttt{Alm\_LoRate} & \texttt{T\_AlmLim} &  & Low flow rate \\
\texttt{Alm\_Overrun} & \texttt{T\_AlmLim} &  & Batch overrun \\
\texttt{Alm\_RateBad} & \texttt{T\_Alarm} &  & Rate signal bad while totalising \\
\bottomrule
\end{longtable}
\end{small}

\subsection{T\_HeatEx}\label{sec:pclib-t-heatex}

Heat exchanger performance monitor: duty, LMTD, UA, fouling, DP.

\par\smallskip{\small \textbf{ISA-88 role:} Control module (monitoring) \quad \textbf{MTP:} - \quad \textbf{Standards:} ISA-88 /\allowbreak{} IEC 61512-1, ISA-18.2 /\allowbreak{} IEC 62682, ISA-101.01}\par

\begin{small}
\begin{longtable}{@{}P{3.70cm}P{2.20cm}P{1.70cm}P{8.91cm}@{}}
\caption{T\_HeatEx Members}\label{tab:pclib-t-heatex} \\
\toprule
\rowcolor{tblhead}\textbf{Member} & \textbf{Type} & \textbf{Init} & \textbf{Description} \\
\midrule
\endfirsthead
\multicolumn{4}{@{}l}{\small\textit{Table \thetable{} continued}} \\
\toprule
\rowcolor{tblhead}\textbf{Member} & \textbf{Type} & \textbf{Init} & \textbf{Description} \\
\midrule
\endhead
\texttt{Cfg\_Tag} & \texttt{STRING[20]} &  & Equipment /\allowbreak{} instrument tag, e.g. FT-101 \\
\texttt{Cfg\_Desc} & \texttt{STRING[40]} &  & Description shown on HMI faceplate \\
\texttt{Inp\_HotInT} & \texttt{REAL} &  & Hot side inlet temperature (degC) \\
\texttt{Inp\_HotOutT} & \texttt{REAL} &  & Hot side outlet temperature (degC) \\
\texttt{Inp\_ColdInT} & \texttt{REAL} &  & Cold side inlet temperature (degC) \\
\texttt{Inp\_ColdOutT} & \texttt{REAL} &  & Cold side outlet temperature (degC) \\
\texttt{Inp\_HotFlow} & \texttt{REAL} &  & Hot side mass flow (kg/\allowbreak{}s) \\
\texttt{Inp\_ColdFlow} & \texttt{REAL} &  & Cold side mass flow (kg/\allowbreak{}s) \\
\texttt{Inp\_HotDP} & \texttt{REAL} &  & Hot side differential pressure (kPa) \\
\texttt{Inp\_ColdDP} & \texttt{REAL} &  & Cold side differential pressure (kPa) \\
\texttt{Cfg\_HotCp} & \texttt{REAL} & 4.186 & Hot side specific heat (kJ/\allowbreak{}kg.K) \\
\texttt{Cfg\_ColdCp} & \texttt{REAL} & 4.186 & Cold side specific heat (kJ/\allowbreak{}kg.K) \\
\texttt{Cfg\_CounterFlow} & \texttt{BOOL} & TRUE & Counter-current (FALSE = co-current) for LMTD \\
\texttt{Cfg\_DesignUA} & \texttt{REAL} & 0.0 & Clean design UA (kW/\allowbreak{}K) \\
\texttt{Cfg\_MinFlow} & \texttt{REAL} & 0.0 & Below this flow the calculations are frozen (kg/\allowbreak{}s) \\
\texttt{Cfg\_CalcFilt} & \texttt{REAL} & 60.0 & Filter time constant for calculated values (s) \\
\texttt{Val\_DutyHot} & \texttt{REAL} &  & Heat released by hot side (kW) \\
\texttt{Val\_DutyCold} & \texttt{REAL} &  & Heat absorbed by cold side (kW) \\
\texttt{Val\_Duty} & \texttt{REAL} &  & Average duty (kW) \\
\texttt{Val\_HeatBal} & \texttt{REAL} &  & Heat balance error (\%) \\
\texttt{Val\_LMTD} & \texttt{REAL} &  & Log mean temperature difference (K) \\
\texttt{Val\_UA} & \texttt{REAL} &  & Actual UA (kW/\allowbreak{}K) \\
\texttt{Val\_Clean} & \texttt{REAL} &  & Cleanliness: UA as \% of design \\
\texttt{Val\_Approach} & \texttt{REAL} &  & Approach temperature (K) \\
\texttt{Val\_Effect} & \texttt{REAL} &  & Thermal effectiveness (0-1) \\
\texttt{Sts\_CalcValid} & \texttt{BOOL} &  & Inputs valid, calculations live \\
\texttt{Alm\_Fouling} & \texttt{T\_AlmLim} &  & Low cleanliness (Val\_Clean below limit) \\
\texttt{Alm\_HotDP} & \texttt{T\_AlmLim} &  & High hot side DP \\
\texttt{Alm\_ColdDP} & \texttt{T\_AlmLim} &  & High cold side DP \\
\texttt{Alm\_HeatBal} & \texttt{T\_AlmLim} &  & Heat balance error (possible instrument fault) \\
\texttt{Alm\_HiOutT} & \texttt{T\_AlmLim} &  & High process outlet temperature \\
\bottomrule
\end{longtable}
\end{small}

\subsection{T\_DosePump}\label{sec:pclib-t-dosepump}

Chemical dosing /\allowbreak{} metering pump: manual, flow-paced, trim, batch.

\par\smallskip{\small \textbf{ISA-88 role:} Equipment module \quad \textbf{MTP:} MonAnaDrv (drive) \quad \textbf{Standards:} ISA-88 /\allowbreak{} IEC 61512-1, ISA-18.2 /\allowbreak{} IEC 62682, ISA-101.01}\par

\begin{small}
\begin{longtable}{@{}P{3.70cm}P{2.20cm}P{1.70cm}P{8.91cm}@{}}
\caption{T\_DosePump Members}\label{tab:pclib-t-dosepump} \\
\toprule
\rowcolor{tblhead}\textbf{Member} & \textbf{Type} & \textbf{Init} & \textbf{Description} \\
\midrule
\endfirsthead
\multicolumn{4}{@{}l}{\small\textit{Table \thetable{} continued}} \\
\toprule
\rowcolor{tblhead}\textbf{Member} & \textbf{Type} & \textbf{Init} & \textbf{Description} \\
\midrule
\endhead
\texttt{Mtr} & \texttt{T\_Motor} &  & Pump motor /\allowbreak{} enable \\
\texttt{Cfg\_RateEU} & \texttt{STRING[8]} &  & Dose rate units text, e.g. L/\allowbreak{}h \\
\texttt{Cfg\_DoseMode} & \texttt{INT} & 0 & 0=Manual rate 1=Flow-paced 2=Flow-paced + trim 3=Batch volume \\
\texttt{Cfg\_MaxRate} & \texttt{REAL} & 0.0 & Pump rate at 100\% output (L/\allowbreak{}h) \\
\texttt{Cfg\_CalFactor} & \texttt{REAL} & 1.0 & Calibration factor from drawdown test \\
\texttt{Cfg\_DoseConc} & \texttt{REAL} & 0.0 & Target dose in process (mg/\allowbreak{}L) \\
\texttt{Cfg\_ChemConc} & \texttt{REAL} & 100.0 & Chemical strength (\% w/\allowbreak{}w) \\
\texttt{Cfg\_ChemSG} & \texttt{REAL} & 1.0 & Chemical specific gravity \\
\texttt{Cfg\_TankCap} & \texttt{REAL} & 0.0 & Chemical day tank capacity (L) \\
\texttt{Cfg\_NoFlowDly} & \texttt{TIME} & T\#30S & Flow verification delay \\
\texttt{Inp\_PaceFlow} & \texttt{REAL} &  & Process flow for pacing (m3/\allowbreak{}h) \\
\texttt{Inp\_TrimPct} & \texttt{REAL} &  & Trim from residual PID (\%, -100..+100) \\
\texttt{Inp\_FlowOK} & \texttt{BOOL} &  & Flow verification switch: dosing detected \\
\texttt{Inp\_ChemLvl} & \texttt{REAL} &  & Chemical day tank level (\%) \\
\texttt{Inp\_Leak} & \texttt{BOOL} &  & Leak detector /\allowbreak{} bund alarm \\
\texttt{Cmd\_Rate} & \texttt{REAL} & 0.0 & Operator dose rate (L/\allowbreak{}h) \\
\texttt{Prog\_Rate} & \texttt{REAL} & 0.0 & Program dose rate (L/\allowbreak{}h) \\
\texttt{Cmd\_CalDone} & \texttt{BOOL} &  & Calibration complete, stamp date (one-shot) \\
\texttt{Out\_Rate} & \texttt{REAL} &  & Speed /\allowbreak{} stroke output to pump (\%) \\
\texttt{Out\_IntlkOK} & \texttt{BOOL} &  & Wire to Mtr.Inp\_IntlkOK \\
\texttt{Val\_RateSP} & \texttt{REAL} &  & Active dose rate (L/\allowbreak{}h) \\
\texttt{Val\_Total} & \texttt{LREAL} &  & Total dosed (L) \\
\texttt{Val\_Today} & \texttt{REAL} &  & Dosed today (L) \\
\texttt{Val\_DaysLeft} & \texttt{REAL} &  & Days of chemical remaining at current usage \\
\texttt{Val\_LastCal} & \texttt{DT} &  & Date of last calibration \\
\texttt{Sts\_RateLimited} & \texttt{BOOL} &  & Demand exceeds pump capacity \\
\texttt{Alm\_NoFlow} & \texttt{T\_Alarm} &  & No flow detected while running \\
\texttt{Alm\_LoChem} & \texttt{T\_AlmLim} &  & Low chemical level (\%) \\
\texttt{Alm\_LoLoChem} & \texttt{T\_AlmLim} &  & Low-low chemical level (\%) --- stops pump \\
\texttt{Alm\_Leak} & \texttt{T\_Alarm} &  & Leak detected \\
\texttt{Alm\_RateLim} & \texttt{T\_Alarm} &  & Dose demand exceeds pump capacity \\
\bottomrule
\end{longtable}
\end{small}

\section{Analyzers and Safety}

\subsection{T\_Analyzer}\label{sec:pclib-t-analyzer}

Process analyzer (pH, ORP, conductivity, DO, turbidity, Cl2 ...).

\par\smallskip{\small \textbf{ISA-88 role:} Control module \quad \textbf{MTP:} AnaMon \quad \textbf{Standards:} ISA-88 /\allowbreak{} IEC 61512-1, ISA-5.1, ISA-18.2 /\allowbreak{} IEC 62682, ISA-101.01, NAMUR NE107, VDI/\allowbreak{}VDE/\allowbreak{}NAMUR 2658-4 (MTP)}\par

\begin{small}
\begin{longtable}{@{}P{3.70cm}P{2.20cm}P{1.70cm}P{8.91cm}@{}}
\caption{T\_Analyzer Members}\label{tab:pclib-t-analyzer} \\
\toprule
\rowcolor{tblhead}\textbf{Member} & \textbf{Type} & \textbf{Init} & \textbf{Description} \\
\midrule
\endfirsthead
\multicolumn{4}{@{}l}{\small\textit{Table \thetable{} continued}} \\
\toprule
\rowcolor{tblhead}\textbf{Member} & \textbf{Type} & \textbf{Init} & \textbf{Description} \\
\midrule
\endhead
\texttt{AI} & \texttt{T\_AI} &  & Measured value, scaling and alarms \\
\texttt{Cfg\_AnType} & \texttt{INT} & 1 & 1=pH 2=ORP 3=Conductivity 4=DO 5=Turbidity 6=Free Cl2 7=TOC 8=Other \\
\texttt{Inp\_Temp} & \texttt{REAL} &  & Sample /\allowbreak{} compensation temperature (degC) \\
\texttt{Inp\_DevFault} & \texttt{BOOL} &  & Transmitter diagnostic fault \\
\texttt{Inp\_DevMaint} & \texttt{BOOL} &  & Transmitter in maintenance /\allowbreak{} calibration \\
\texttt{Inp\_DevWarn} & \texttt{BOOL} &  & Transmitter warning (e.g. sensor ageing) \\
\texttt{Inp\_NE107} & \texttt{INT} & 0 & NE107 status reported by smart device (HART /\allowbreak{} fieldbus) \\
\texttt{Sts\_NE107} & \texttt{INT} &  & NAMUR NE107 status: 0=Good 1=Failure (F) 2=Function check (C) 3=Out of specification (S) 4=Maintenance required (M) \\
\texttt{Inp\_SampleOK} & \texttt{BOOL} &  & Sample flow healthy (flow switch) \\
\texttt{Inp\_Slope} & \texttt{REAL} &  & Sensor slope from transmitter (e.g. pH mV/\allowbreak{}pH) \\
\texttt{Cfg\_HoldMax} & \texttt{TIME} & T\#1H & Maximum hold time before alarm \\
\texttt{Cfg\_CalDays} & \texttt{INT} & 30 & Calibration interval (days) \\
\texttt{Cfg\_FrozenTime} & \texttt{TIME} & T\#30M & PV unchanged for this long = frozen sensor \\
\texttt{Cfg\_FrozenBand} & \texttt{REAL} & 0.0 & Band within which PV is considered unchanged \\
\texttt{Cmd\_Hold} & \texttt{BOOL} &  & Hold PV and suppress alarms for calibration (one-shot) \\
\texttt{Cmd\_Release} & \texttt{BOOL} &  & Release hold (one-shot) \\
\texttt{Cmd\_CalDone} & \texttt{BOOL} &  & Calibration completed, stamp date (one-shot) \\
\texttt{Val\_PV} & \texttt{REAL} &  & Value for control --- held value while in hold \\
\texttt{Val\_HeldPV} & \texttt{REAL} &  & Value captured at start of hold \\
\texttt{Val\_LastCal} & \texttt{DT} &  & Date/\allowbreak{}time of last calibration \\
\texttt{Val\_DaysSinceCal} & \texttt{REAL} &  & Days since last calibration \\
\texttt{Sts\_Hold} & \texttt{BOOL} &  & In hold /\allowbreak{} calibration \\
\texttt{Sts\_CalDue} & \texttt{BOOL} &  & Calibration due \\
\texttt{Alm\_DevFault} & \texttt{T\_Alarm} &  & Transmitter fault \\
\texttt{Alm\_NoSample} & \texttt{T\_Alarm} &  & No sample flow \\
\texttt{Alm\_Frozen} & \texttt{T\_Alarm} &  & Frozen /\allowbreak{} stuck reading \\
\texttt{Alm\_CalDue} & \texttt{T\_Alarm} &  & Calibration overdue \\
\texttt{Alm\_HoldLong} & \texttt{T\_Alarm} &  & Hold exceeded Cfg\_HoldMax \\
\bottomrule
\end{longtable}
\end{small}

\subsection{T\_GasDet}\label{sec:pclib-t-gasdet}

Fixed gas detector (4-20mA): latching A1/\allowbreak{}A2/\allowbreak{}A3, inhibit, fault.

\par\smallskip{\small \textbf{ISA-88 role:} Control module \quad \textbf{MTP:} AnaMon \quad \textbf{Standards:} ISA-88 /\allowbreak{} IEC 61512-1, ISA-5.1, ISA-18.2 /\allowbreak{} IEC 62682, ISA-101.01, NAMUR NE107, VDI/\allowbreak{}VDE/\allowbreak{}NAMUR 2658-4 (MTP), NAMUR NE43, ISA-84 /\allowbreak{} IEC 61511 (BPCS layer only)}\par

\begin{small}
\begin{longtable}{@{}P{3.70cm}P{2.20cm}P{1.70cm}P{8.91cm}@{}}
\caption{T\_GasDet Members}\label{tab:pclib-t-gasdet} \\
\toprule
\rowcolor{tblhead}\textbf{Member} & \textbf{Type} & \textbf{Init} & \textbf{Description} \\
\midrule
\endfirsthead
\multicolumn{4}{@{}l}{\small\textit{Table \thetable{} continued}} \\
\toprule
\rowcolor{tblhead}\textbf{Member} & \textbf{Type} & \textbf{Init} & \textbf{Description} \\
\midrule
\endhead
\texttt{Cfg\_Tag} & \texttt{STRING[20]} &  & Equipment /\allowbreak{} instrument tag, e.g. FT-101 \\
\texttt{Cfg\_Desc} & \texttt{STRING[40]} &  & Description shown on HMI faceplate \\
\texttt{Cfg\_Gas} & \texttt{STRING[12]} &  & Gas, e.g. CH4, H2S, CO, O2 \\
\texttt{Cfg\_EU} & \texttt{STRING[8]} &  & Units, e.g. \%LEL, ppm, \%vol \\
\texttt{Inp\_Raw} & \texttt{INT} &  & Raw counts from analog channel \\
\texttt{Inp\_ChFault} & \texttt{BOOL} &  & Channel /\allowbreak{} module diagnostic fault \\
\texttt{Cfg\_RawMin} & \texttt{INT} & 0 & Raw counts at 4 mA \\
\texttt{Cfg\_RawMax} & \texttt{INT} & 10000 & Raw counts at 20 mA \\
\texttt{Cfg\_RangeMax} & \texttt{REAL} & 100.0 & Concentration at 20 mA \\
\texttt{Cfg\_O2Mode} & \texttt{BOOL} &  & Oxygen depletion: A1/\allowbreak{}A2 act on falling concentration \\
\texttt{Cfg\_FaultmA} & \texttt{REAL} & 3.0 & Signal below this (mA) = detector fault \\
\texttt{Cfg\_MaintLomA} & \texttt{REAL} & 0.0 & Maintenance /\allowbreak{} inhibit signal band low (mA), vendor-specific \\
\texttt{Cfg\_MaintHimA} & \texttt{REAL} & 0.0 & Maintenance /\allowbreak{} inhibit signal band high (mA), vendor-specific \\
\texttt{Cfg\_OverRngmA} & \texttt{REAL} & 20.5 & Signal above this (mA) = over-range (treat as alarm) \\
\texttt{Cfg\_InhibitMax} & \texttt{TIME} & T\#2H & Maximum operator inhibit time \\
\texttt{Cmd\_Inhibit} & \texttt{BOOL} &  & Operator inhibit for maintenance (one-shot) \\
\texttt{Cmd\_Uninhibit} & \texttt{BOOL} &  & Remove inhibit (one-shot) \\
\texttt{Cmd\_Reset} & \texttt{BOOL} &  & Reset latched alarms (one-shot) \\
\texttt{Sim\_En} & \texttt{BOOL} &  & Simulation enabled \\
\texttt{Sim\_mA} & \texttt{REAL} & 4.0 & Simulated signal (mA) \\
\texttt{Out\_A1} & \texttt{BOOL} &  & Alarm level 1 active (to voting /\allowbreak{} beacons) \\
\texttt{Out\_A2} & \texttt{BOOL} &  & Alarm level 2 active \\
\texttt{Out\_A3} & \texttt{BOOL} &  & Alarm level 3 active \\
\texttt{Out\_Fault} & \texttt{BOOL} &  & Detector fault /\allowbreak{} unavailable \\
\texttt{Val\_mA} & \texttt{REAL} &  & Signal (mA) \\
\texttt{Val\_Conc} & \texttt{REAL} &  & Gas concentration (Cfg\_EU) \\
\texttt{Val\_Peak} & \texttt{REAL} &  & Peak concentration since last reset \\
\texttt{Sts\_Inhibit} & \texttt{BOOL} &  & Operator inhibit active \\
\texttt{Sts\_DetMaint} & \texttt{BOOL} &  & Detector signalling maintenance \\
\texttt{Sts\_InhibitLeft} & \texttt{TIME} &  & Inhibit time remaining \\
\texttt{Sts\_Sim} & \texttt{BOOL} &  & Simulation active \\
\texttt{Inp\_NE107} & \texttt{INT} & 0 & NE107 status reported by smart device (HART /\allowbreak{} fieldbus) \\
\texttt{Sts\_NE107} & \texttt{INT} &  & NAMUR NE107 status: 0=Good 1=Failure (F) 2=Function check (C) 3=Out of specification (S) 4=Maintenance required (M) \\
\texttt{Alm\_A1} & \texttt{T\_AlmLim} &  & Alarm level 1 (low) --- set Cfg\_Latch as required \\
\texttt{Alm\_A2} & \texttt{T\_AlmLim} &  & Alarm level 2 (high) --- set Cfg\_Latch as required \\
\texttt{Alm\_A3} & \texttt{T\_AlmLim} &  & Alarm level 3 (high-high /\allowbreak{} STEL) \\
\texttt{Alm\_Fault} & \texttt{T\_Alarm} &  & Detector fault \\
\texttt{Alm\_OverRng} & \texttt{T\_Alarm} &  & Over-range \\
\texttt{Alm\_InhibitLong} & \texttt{T\_Alarm} &  & Inhibit timed out \\
\bottomrule
\end{longtable}
\end{small}

\subsection{T\_Intlk}\label{sec:pclib-t-intlk}

Interlock /\allowbreak{} permissive group: 16 conditions, first-out, timed bypass.

\par\smallskip{\small \textbf{ISA-88 role:} Support (interlock group) \quad \textbf{MTP:} LockView16 \quad \textbf{Standards:} ISA-84 /\allowbreak{} IEC 61511 (BPCS layer only), ISA-18.2 /\allowbreak{} IEC 62682, ISA-101.01, VDI/\allowbreak{}VDE/\allowbreak{}NAMUR 2658-4 (MTP)}\par

\begin{small}
\begin{longtable}{@{}P{3.70cm}P{2.20cm}P{1.70cm}P{8.91cm}@{}}
\caption{T\_Intlk Members}\label{tab:pclib-t-intlk} \\
\toprule
\rowcolor{tblhead}\textbf{Member} & \textbf{Type} & \textbf{Init} & \textbf{Description} \\
\midrule
\endfirsthead
\multicolumn{4}{@{}l}{\small\textit{Table \thetable{} continued}} \\
\toprule
\rowcolor{tblhead}\textbf{Member} & \textbf{Type} & \textbf{Init} & \textbf{Description} \\
\midrule
\endhead
\texttt{Cfg\_Tag} & \texttt{STRING[20]} &  & Equipment /\allowbreak{} instrument tag, e.g. FT-101 \\
\texttt{Cfg\_Desc} & \texttt{STRING[40]} &  & Description shown on HMI faceplate \\
\texttt{Inp\_Cond} & \texttt{ARRAY[0..15] OF BOOL} &  & Condition n: TRUE = healthy /\allowbreak{} permitted \\
\texttt{Cfg\_Used} & \texttt{ARRAY[0..15] OF BOOL} &  & Condition n is configured \\
\texttt{Cfg\_CondDesc} & \texttt{ARRAY[0..15] OF STRING[32]} &  & HMI text for condition n \\
\texttt{Cfg\_Bypassable} & \texttt{ARRAY[0..15] OF BOOL} &  & Condition n may be bypassed \\
\texttt{Cfg\_Latch} & \texttt{BOOL} &  & Latch failures until Cmd\_Reset (interlock); FALSE = permissive \\
\texttt{Cfg\_BypassMax} & \texttt{TIME} & T\#8H & Maximum bypass duration \\
\texttt{Cmd\_Bypass} & \texttt{ARRAY[0..15] OF BOOL} &  & Bypass condition n (one-shot) \\
\texttt{Cmd\_Unbypass} & \texttt{ARRAY[0..15] OF BOOL} &  & Remove bypass on condition n (one-shot) \\
\texttt{Cmd\_UnbypassAll} & \texttt{BOOL} &  & Remove all bypasses (one-shot) \\
\texttt{Inp\_BypassPerm} & \texttt{BOOL} &  & Bypass authorisation (key switch /\allowbreak{} HMI permit) --- IEC 61511 bypass management \\
\texttt{Cmd\_Reset} & \texttt{BOOL} &  & Reset latched conditions (one-shot) \\
\texttt{Out\_OK} & \texttt{BOOL} &  & All used conditions healthy or bypassed --- wire to Inp\_IntlkOK /\allowbreak{} Inp\_PermOK \\
\texttt{Sts\_OKNoByp} & \texttt{BOOL} &  & All used conditions healthy without bypasses \\
\texttt{Sts\_Bypassed} & \texttt{ARRAY[0..15] OF BOOL} &  & Condition n bypassed \\
\texttt{Sts\_Latched} & \texttt{ARRAY[0..15] OF BOOL} &  & Condition n latched failed \\
\texttt{Sts\_FirstOut} & \texttt{INT} &  & First condition to fail (-1 = none) \\
\texttt{Sts\_NumBypassed} & \texttt{INT} &  & Number of bypassed conditions \\
\texttt{Sts\_BypassLeft} & \texttt{TIME} &  & Time left on the oldest bypass \\
\texttt{Sts\_ResetReq} & \texttt{BOOL} &  & Conditions healthy, reset required \\
\texttt{Sts\_CondWord} & \texttt{WORD} &  & Packed Inp\_Cond (bit n = condition n) for SCADA \\
\texttt{Sts\_BypWord} & \texttt{WORD} &  & Packed Sts\_Bypassed for SCADA \\
\texttt{Sts\_LatchWord} & \texttt{WORD} &  & Packed Sts\_Latched for SCADA \\
\texttt{Alm\_Trip} & \texttt{T\_Alarm} &  & Interlock tripped \\
\texttt{Alm\_Bypass} & \texttt{T\_Alarm} &  & One or more conditions bypassed \\
\bottomrule
\end{longtable}
\end{small}

\subsection{T\_ESD}\label{sec:pclib-t-esd}

Emergency shutdown group (BPCS layer): de-energise-to-trip, MOS, first-out.

\par\smallskip{\small \textbf{ISA-88 role:} Support (shutdown group) \quad \textbf{MTP:} - \quad \textbf{Standards:} ISA-84 /\allowbreak{} IEC 61511 (BPCS layer only), ISA-18.2 /\allowbreak{} IEC 62682, ISA-101.01}\par

\begin{small}
\begin{longtable}{@{}P{3.70cm}P{2.20cm}P{1.70cm}P{8.91cm}@{}}
\caption{T\_ESD Members}\label{tab:pclib-t-esd} \\
\toprule
\rowcolor{tblhead}\textbf{Member} & \textbf{Type} & \textbf{Init} & \textbf{Description} \\
\midrule
\endfirsthead
\multicolumn{4}{@{}l}{\small\textit{Table \thetable{} continued}} \\
\toprule
\rowcolor{tblhead}\textbf{Member} & \textbf{Type} & \textbf{Init} & \textbf{Description} \\
\midrule
\endhead
\texttt{Cfg\_Tag} & \texttt{STRING[20]} &  & Equipment /\allowbreak{} instrument tag, e.g. FT-101 \\
\texttt{Cfg\_Desc} & \texttt{STRING[40]} &  & Description shown on HMI faceplate \\
\texttt{Cfg\_Level} & \texttt{INT} & 1 & ESD level, e.g. 0=Plant 1=Unit 2=Equipment \\
\texttt{Inp\_OK} & \texttt{ARRAY[0..15] OF BOOL} &  & Trip input n: TRUE = healthy \\
\texttt{Inp\_UpstrOK} & \texttt{BOOL} & TRUE & Higher-level ESD healthy (cascade) \\
\texttt{Inp\_ManTripOK} & \texttt{BOOL} &  & Field ESD push-button circuit healthy \\
\texttt{Inp\_ResetPB} & \texttt{BOOL} &  & Field reset push-button \\
\texttt{Inp\_MOSKey} & \texttt{BOOL} &  & Maintenance override key switch enabled \\
\texttt{Inp\_FinalSafe} & \texttt{BOOL} &  & All final elements confirmed in safe state \\
\texttt{Cfg\_Used} & \texttt{ARRAY[0..15] OF BOOL} &  & Input n configured \\
\texttt{Cfg\_InDesc} & \texttt{ARRAY[0..15] OF STRING[32]} &  & HMI text for input n \\
\texttt{Cfg\_MOSAllowed} & \texttt{ARRAY[0..15] OF BOOL} &  & Override allowed on input n \\
\texttt{Cfg\_MOSMax} & \texttt{TIME} & T\#4H & Maximum override duration \\
\texttt{Cfg\_FieldReset} & \texttt{BOOL} &  & Field reset push-button required to reset \\
\texttt{Cfg\_SafeTime} & \texttt{TIME} & T\#30S & Time for final elements to reach safe state \\
\texttt{Cmd\_ManTrip} & \texttt{BOOL} &  & HMI manual trip (one-shot) \\
\texttt{Cmd\_Reset} & \texttt{BOOL} &  & HMI reset (one-shot) \\
\texttt{Cmd\_MOS} & \texttt{ARRAY[0..15] OF BOOL} &  & Apply override on input n (one-shot) \\
\texttt{Cmd\_MOSOff} & \texttt{ARRAY[0..15] OF BOOL} &  & Remove override on input n (one-shot) \\
\texttt{Out\_Healthy} & \texttt{BOOL} &  & TRUE = run permitted (energised); FALSE = trip \\
\texttt{Sts\_Tripped} & \texttt{BOOL} &  & Group tripped \\
\texttt{Sts\_ResetReady} & \texttt{BOOL} &  & All inputs healthy, reset possible \\
\texttt{Sts\_FirstOut} & \texttt{INT} &  & First input to trip (-1 none, 16=manual, 17=upstream) \\
\texttt{Sts\_MOS} & \texttt{ARRAY[0..15] OF BOOL} &  & Override active on input n \\
\texttt{Sts\_NumMOS} & \texttt{INT} &  & Number of active overrides \\
\texttt{Sts\_InWord} & \texttt{WORD} &  & Packed Inp\_OK for SCADA \\
\texttt{Sts\_LatchWord} & \texttt{WORD} &  & Packed latched trips for SCADA \\
\texttt{Sts\_MOSWord} & \texttt{WORD} &  & Packed Sts\_MOS for SCADA \\
\texttt{Val\_TripTime} & \texttt{DT} &  & Date/\allowbreak{}time of last trip \\
\texttt{Val\_TripCount} & \texttt{UDINT} &  & Trip counter \\
\texttt{Alm\_Trip} & \texttt{T\_Alarm} &  & ESD trip \\
\texttt{Alm\_FailSafe} & \texttt{T\_Alarm} &  & Final elements failed to reach safe state \\
\texttt{Alm\_MOS} & \texttt{T\_Alarm} &  & Maintenance override active \\
\texttt{Alm\_MOSLong} & \texttt{T\_Alarm} &  & Override exceeded Cfg\_MOSMax \\
\bottomrule
\end{longtable}
\end{small}

\section{ISA-88 Batch and Procedural}

\subsection{T\_Phase}\label{sec:pclib-t-phase}

ISA-88 equipment phase: Part 1 state model, commands, modes, parameters, reports.

\par\smallskip{\small \textbf{ISA-88 role:} Equipment phase (procedural) \quad \textbf{MTP:} Service (state machine) \quad \textbf{Standards:} ISA-88 /\allowbreak{} IEC 61512-1 Part 1 state model, VDI/\allowbreak{}VDE/\allowbreak{}NAMUR 2658-4 (MTP)}\par

\begin{small}
\begin{longtable}{@{}P{3.70cm}P{2.20cm}P{1.70cm}P{8.91cm}@{}}
\caption{T\_Phase Members}\label{tab:pclib-t-phase} \\
\toprule
\rowcolor{tblhead}\textbf{Member} & \textbf{Type} & \textbf{Init} & \textbf{Description} \\
\midrule
\endfirsthead
\multicolumn{4}{@{}l}{\small\textit{Table \thetable{} continued}} \\
\toprule
\rowcolor{tblhead}\textbf{Member} & \textbf{Type} & \textbf{Init} & \textbf{Description} \\
\midrule
\endhead
\texttt{Cfg\_Tag} & \texttt{STRING[20]} &  & Equipment /\allowbreak{} instrument tag, e.g. FT-101 \\
\texttt{Cfg\_Desc} & \texttt{STRING[40]} &  & Description shown on HMI faceplate \\
\texttt{Cfg\_PhaseID} & \texttt{INT} &  & Phase ID (unique in the unit) for batch server /\allowbreak{} unit procedure \\
\texttt{Cfg\_UnitID} & \texttt{INT} &  & ID of the unit this phase belongs to \\
\texttt{Cfg\_StepTimeout} & \texttt{TIME} & T\#0S & Step watchdog, T\#0S = off \\
\texttt{Cmd\_Auto} & \texttt{BOOL} &  & Operator request Automatic mode (one-shot) \\
\texttt{Cmd\_SemiAuto} & \texttt{BOOL} &  & Operator request Semi-automatic mode (one-shot) \\
\texttt{Cmd\_Man} & \texttt{BOOL} &  & Operator request Manual mode (one-shot) \\
\texttt{Cmd\_Start} & \texttt{BOOL} &  & Operator Start (one-shot) \\
\texttt{Cmd\_Stop} & \texttt{BOOL} &  & Operator Stop (one-shot) \\
\texttt{Cmd\_Hold} & \texttt{BOOL} &  & Operator Hold (one-shot) \\
\texttt{Cmd\_Restart} & \texttt{BOOL} &  & Operator Restart from Held (one-shot) \\
\texttt{Cmd\_Abort} & \texttt{BOOL} &  & Operator Abort (one-shot) \\
\texttt{Cmd\_Reset} & \texttt{BOOL} &  & Operator Reset to Idle (one-shot) \\
\texttt{Cmd\_Pause} & \texttt{BOOL} &  & Operator Pause at next transition (one-shot) \\
\texttt{Cmd\_Resume} & \texttt{BOOL} &  & Operator Resume from Paused (one-shot) \\
\texttt{Cmd\_Advance} & \texttt{BOOL} &  & Advance to next step in Semi-automatic (one-shot) \\
\texttt{Cmd\_TakeOwn} & \texttt{BOOL} &  & Operator takes control from the recipe (one-shot) \\
\texttt{Prog\_Cmd} & \texttt{INT} &  & Command from batch server /\allowbreak{} unit procedure: 1=Start 2=Stop 3=Hold 4=Restart 5=Abort 6=Reset 7=Pause 8=Resume \\
\texttt{Prog\_Acquire} & \texttt{BOOL} &  & Recipe (batch server /\allowbreak{} unit procedure) takes ownership \\
\texttt{Prog\_Release} & \texttt{BOOL} &  & Recipe releases ownership \\
\texttt{Prog\_ReqAck} & \texttt{BOOL} &  & Batch server acknowledges Sts\_ReqCode \\
\texttt{Prm\_Real} & \texttt{ARRAY[1..16] OF REAL} &  & Recipe parameters (REAL) --- control recipe formula \\
\texttt{Prm\_DInt} & \texttt{ARRAY[1..8] OF DINT} &  & Recipe parameters (DINT) --- control recipe formula \\
\texttt{Rpt\_Real} & \texttt{ARRAY[1..16] OF REAL} &  & Report values (REAL) returned to batch record \\
\texttt{Rpt\_DInt} & \texttt{ARRAY[1..8] OF DINT} &  & Report values (DINT) returned to batch record \\
\texttt{Val\_BatchID} & \texttt{STRING[32]} &  & Batch ID of the control recipe running this phase \\
\texttt{Val\_StateTime} & \texttt{TIME} &  & Time in current state \\
\texttt{Val\_RunTime} & \texttt{TIME} &  & Running time excluding Held /\allowbreak{} Paused \\
\texttt{Val\_FailText} & \texttt{STRING[40]} &  & Failure description for HMI /\allowbreak{} batch record \\
\texttt{Sts\_State} & \texttt{INT} &  & ISA-88 state: 0=Idle 1=Running 2=Complete 3=Pausing 4=Paused 5=Holding 6=Held 7=Restarting 8=Stopping 9=Stopped 10=Aborting 11=Aborted \\
\texttt{Sts\_PrevState} & \texttt{INT} &  & Previous state (for Restart /\allowbreak{} diagnostics) \\
\texttt{Sts\_Mode} & \texttt{INT} &  & ISA-88 mode: 1=Manual 2=Semi-automatic 3=Automatic \\
\texttt{Sts\_Owner} & \texttt{INT} &  & 0=Operator 1=Recipe (batch server /\allowbreak{} unit procedure) \\
\texttt{Sts\_CmdMask} & \texttt{WORD} &  & Commands valid in current state (bit n = command n of Prog\_Cmd) --- for HMI \\
\texttt{Sts\_CmdRejected} & \texttt{BOOL} &  & Last command was not valid in the current state \\
\texttt{Sts\_Step} & \texttt{INT} &  & Active step number within the state logic \\
\texttt{Sts\_WaitAdvance} & \texttt{BOOL} &  & Semi-automatic: waiting for Cmd\_Advance \\
\texttt{Sts\_ReqCode} & \texttt{INT} &  & Request to batch server (e.g. download params, upload reports), 0 = none \\
\texttt{Sts\_FailCode} & \texttt{INT} &  & Failure code, 0 = none \\
\texttt{Alm\_Fail} & \texttt{T\_Alarm} &  & Phase failure \\
\texttt{Alm\_StepTimeout} & \texttt{T\_Alarm} &  & Step exceeded Cfg\_StepTimeout \\
\bottomrule
\end{longtable}
\end{small}

\subsection{T\_Unit}\label{sec:pclib-t-unit}

ISA-88 unit: allocation, batch identity, procedural and equipment state.

\par\smallskip{\small \textbf{ISA-88 role:} Unit (physical) \quad \textbf{MTP:} - \quad \textbf{Standards:} ISA-88 /\allowbreak{} IEC 61512-1, ISA-18.2 /\allowbreak{} IEC 62682, ISA-101.01}\par

\begin{small}
\begin{longtable}{@{}P{3.70cm}P{2.20cm}P{1.70cm}P{8.91cm}@{}}
\caption{T\_Unit Members}\label{tab:pclib-t-unit} \\
\toprule
\rowcolor{tblhead}\textbf{Member} & \textbf{Type} & \textbf{Init} & \textbf{Description} \\
\midrule
\endfirsthead
\multicolumn{4}{@{}l}{\small\textit{Table \thetable{} continued}} \\
\toprule
\rowcolor{tblhead}\textbf{Member} & \textbf{Type} & \textbf{Init} & \textbf{Description} \\
\midrule
\endhead
\texttt{Cfg\_Tag} & \texttt{STRING[20]} &  & Equipment /\allowbreak{} instrument tag, e.g. FT-101 \\
\texttt{Cfg\_Desc} & \texttt{STRING[40]} &  & Description shown on HMI faceplate \\
\texttt{Cfg\_UnitID} & \texttt{INT} &  & Unit ID \\
\texttt{Cfg\_UnitClass} & \texttt{STRING[24]} &  & Unit class name (master recipe equipment requirement) \\
\texttt{Cfg\_CleanHold} & \texttt{TIME} & T\#72H & Clean hold time --- clean status expires after this \\
\texttt{Inp\_ReqOwnerID} & \texttt{DINT} &  & ID of control recipe requesting allocation \\
\texttt{Prog\_Acquire} & \texttt{BOOL} &  & Allocate unit to requesting control recipe \\
\texttt{Prog\_Release} & \texttt{BOOL} &  & Release unit \\
\texttt{Cmd\_ForceRelease} & \texttt{BOOL} &  & Supervisor forces release (one-shot) \\
\texttt{Prog\_EqState} & \texttt{INT} &  & Equipment state request from phase logic \\
\texttt{Cmd\_EqState} & \texttt{INT} &  & Operator equipment state change (e.g. to Maintenance) \\
\texttt{Val\_OwnerID} & \texttt{DINT} &  & Allocated control recipe ID, 0 = free \\
\texttt{Val\_BatchID} & \texttt{STRING[32]} &  & Batch ID \\
\texttt{Val\_RecipeName} & \texttt{STRING[32]} &  & Control recipe /\allowbreak{} master recipe name \\
\texttt{Val\_RecipeVer} & \texttt{INT} &  & Recipe version \\
\texttt{Val\_ProductCode} & \texttt{STRING[24]} &  & Product /\allowbreak{} material code \\
\texttt{Val\_BatchSize} & \texttt{REAL} &  & Batch size (scaled) \\
\texttt{Val\_BatchStart} & \texttt{DT} &  & Batch start date/\allowbreak{}time \\
\texttt{Val\_LastClean} & \texttt{DT} &  & Date/\allowbreak{}time of last completed clean \\
\texttt{Sts\_Allocated} & \texttt{BOOL} &  & Unit allocated to a control recipe \\
\texttt{Sts\_ProcState} & \texttt{INT} &  & Unit procedure state: 0=Idle 1=Running 2=Complete 3=Pausing 4=Paused 5=Holding 6=Held 7=Restarting 8=Stopping 9=Stopped 10=Aborting 11=Aborted \\
\texttt{Sts\_Mode} & \texttt{INT} &  & Unit procedure mode: 1=Manual 2=Semi-automatic 3=Automatic \\
\texttt{Sts\_ActivePhase} & \texttt{INT} &  & Cfg\_PhaseID of the active phase, 0 = none \\
\texttt{Sts\_EqState} & \texttt{INT} &  & ISA-88 equipment state: 0=Available-clean 1=In use 2=Dirty 3=Cleaning (CIP) 4=Sterilising (SIP) 5=Maintenance 6=Out of service \\
\texttt{Sts\_CleanExpired} & \texttt{BOOL} &  & Clean hold time exceeded \\
\texttt{Alm\_CleanExpired} & \texttt{T\_Alarm} &  & Clean status expired \\
\bottomrule
\end{longtable}
\end{small}

\section{Reference Implementation: FB\_AI}\label{sec:pclib-fbai}
\texttt{FB\_AI} is the first PCLib DFB: an IEC 61131-3 Structured Text block that processes one \texttt{T\_AI} instance. It is generated by \texttt{gen\_dfb.py} and imported from \texttt{export/PCLib\_DFB\_AI\_CE\_V16.2.xdb}, which also contains the shared alarm blocks \texttt{FB\_AlmCore}, \texttt{FB\_Alarm}, and \texttt{FB\_AlmLim}. The same alarm blocks are intended to be reused by every other PCLib DFB as it is written. The block uses no located variables, global variables, or system bits, so all data passes through its parameters. Run it once per scan in a cyclic task; it measures its own elapsed time.

\begin{small}
\begin{verbatim}
(* Variables: FT101 : T_AI;  fbFT101 : FB_AI; *)
FT101.Inp_Raw     := %IW0.2.0;            (* or the M580 channel value *)
FT101.Inp_ChFault := NOT %I0.2.0.ERR;     (* channel health *)
fbFT101(AI := FT101);
\end{verbatim}
\end{small}


\begin{small}
\begin{longtable}{@{}P{0.70cm}P{9.50cm}P{6.73cm}@{}}
\caption{FB\_AI Processing Order}\label{tab:pclib-fbai-order} \\
\toprule
\rowcolor{tblhead}\textbf{\#} & \textbf{Step} & \textbf{Members} \\
\midrule
\endfirsthead
\multicolumn{3}{@{}l}{\small\textit{Table \thetable{} continued}} \\
\toprule
\rowcolor{tblhead}\textbf{\#} & \textbf{Step} & \textbf{Members} \\
\midrule
\endhead
1 & Elapsed time since last scan (free-running TON, no system calls) & - \\
2 & Configuration check: raw or EU span of zero makes the PV bad & \texttt{Cfg\_RawMin/Max}, \texttt{Cfg\_EUMin/Max} \\
3 & Raw counts to \% of span & \texttt{Inp\_Raw} \\
4 & Signal validation: NAMUR NE43 limits, channel fault, device NE107 Failure & \texttt{Cfg\_BadLoPct/HiPct}, \texttt{Inp\_ChFault}, \texttt{Inp\_NE107}, \texttt{Sts\_URng/ORng} \\
5 & Square-root extraction, low cut-off, scaling to EU & \texttt{Cfg\_SqrtExtract}, \texttt{Cfg\_LoCutoff} \\
6 & Simulation replaces the field signal, and field faults are ignored & \texttt{Sim\_En}, \texttt{Sim\_PV}, \texttt{Sts\_Sim} \\
7 & First-order filter. Re-initialised on start-up, bad-to-good recovery and sim change, so there is no bump & \texttt{Cfg\_FiltTC} \\
8 & Bad-value action: hold last good /\allowbreak{} substitute /\allowbreak{} EUMin /\allowbreak{} EUMax & \texttt{Cfg\_BadAction}, \texttt{Cfg\_SubstPV}, \texttt{Sts\_Subst} \\
9 & Quality and NAMUR NE107 status. Simulation reports as C (function check) & \texttt{Sts\_Quality}, \texttt{Sts\_NE107} \\
10 & Rate of change, 1 s window. Re-initialised on recovery and sim change, so there are no false ROC alarms & \texttt{Val\_ROC} \\
11 & Alarms: HH/\allowbreak{}H/\allowbreak{}L/\allowbreak{}LL/\allowbreak{}ROC limit alarms and the bad-PV alarm & \texttt{Alm\_*} \\
\bottomrule
\end{longtable}
\end{small}

Limit alarms are disabled by default (\texttt{Cfg\_En = FALSE}); enable each one after its limit has been set in alarm rationalisation. The simulator test checklist for \texttt{FB\_AI} is kept with the library in the EFC Repo.
